\documentclass[UTF8,fontset=none,zihao=5,a4paper]{ctexart}
\usepackage[margin=22mm,headheight=16pt,headsep=7mm,footskip=10mm]{geometry}
\usepackage{fontspec}
\setmainfont{NimbusRoman}[Path=/usr/share/fonts/opentype/urw-base35/,Extension=.otf,UprightFont=*-Regular,BoldFont=*-Bold,ItalicFont=*-Italic,BoldItalicFont=*-BoldItalic]
\setsansfont{NimbusSans}[Path=/usr/share/fonts/opentype/urw-base35/,Extension=.otf,UprightFont=*-Regular,BoldFont=*-Bold,ItalicFont=*-Italic,BoldItalicFont=*-BoldItalic]
\setmonofont{DejaVu Sans Mono}[Scale=0.80]
\setCJKmainfont[AutoFakeSlant=true]{Noto Serif CJK SC}
\setCJKsansfont{Noto Sans CJK SC}
\setCJKmonofont{Noto Sans Mono CJK SC}
\usepackage{amsmath,amssymb,amsthm,bm,mathtools}
\usepackage{booktabs,longtable,tabularx,array,graphicx,xcolor}
\usepackage{tikz}
\usetikzlibrary{arrows.meta,positioning}
\usepackage{enumitem,fancyhdr,microtype}
\usepackage[colorlinks=true,linkcolor=blue!45!black,citecolor=blue!45!black,urlcolor=blue!55!black,breaklinks=true]{hyperref}
\usepackage{url}
\urlstyle{same}
\definecolor{accent}{HTML}{244F6B}
\definecolor{lightaccent}{HTML}{EDF3F7}
\ctexset{section={format=\Large\bfseries\color{accent}},subsection={format=\large\bfseries},subsubsection={format=\normalsize\bfseries}}
\setlength{\parskip}{3pt}
\linespread{1.12}
\setlength{\emergencystretch}{3em}
\setlist{nosep,leftmargin=2em}
\renewcommand{\arraystretch}{1.17}
\numberwithin{equation}{section}
\newtheorem{proposition}{命题}[section]
\renewcommand{\proofname}{证明}
\DeclareMathOperator{\diag}{diag}
\DeclareMathOperator{\dist}{dist}
\DeclareMathOperator{\clip}{clip}
\DeclareMathOperator{\Cov}{Cov}
\DeclareMathOperator{\Var}{Var}
\DeclareMathOperator{\CVaR}{CVaR}
\DeclareMathOperator*{\argmin}{arg\,min}
\DeclareMathOperator*{\argmax}{arg\,max}
\newcommand{\R}{\mathbb R}
\newcommand{\E}{\mathbb E}
\newcommand{\Prob}{\mathbb P}
\newcommand{\norm}[1]{\left\lVert#1\right\rVert}
\newcommand{\abs}[1]{\left\lvert#1\right\rvert}
\newcommand{\wrap}{\operatorname{wrap}}
\newcommand{\ind}{\mathbf1}
\newcommand{\code}[1]{\texttt{\small\detokenize{#1}}}
\newcommand{\base}[1]{\textbf{原稿依据：}\code{#1}。}
\pagestyle{fancy}
\fancyhf{}
\fancyhead[L]{\small duojin02：固定时延下的 CUDA 算法重构分析}
\fancyhead[R]{\small 技术研究报告}
\fancyfoot[C]{\thepage}
\hypersetup{pdftitle={duojin02 CUDA 加速与固定时延算法重构分析},pdfauthor={技术分析报告},pdfsubject={数学推导、计算预算与验证方案；不修改项目代码}}
\begin{document}
\begin{titlepage}
\vspace*{15mm}
{\color{accent}\Huge\bfseries duojin02\\[4mm]CUDA 加速与算法重构分析\par}
\vspace{9mm}
{\LARGE 在相同在线耗时内，\\利用更大的计算与存储预算改善算法效果\par}
\vspace{12mm}
{\large 基于所提供算法论文源文件的逐项审阅与数学推导\par}
\vspace{7mm}
{\large 2026 年 10 月 15 日\par}
\vspace{14mm}
\noindent\fcolorbox{accent}{lightaccent}{\parbox{0.91\textwidth}{
\textbf{核心判断}\par
最有价值的方向并非把所有 CPU 运算原样搬到 GPU，而是重构计算结构：
用多场景采样控制处理不确定性，用多假设定位保持歧义，
用预计算配置空间与距离场把昂贵的几何运算转成密集查表。
GPU 的并行与显存用于换取更充分的决策，同时用完整时延预算约束规模。\par
这些方向具有可推导的收益机制，但“同耗时且效果更好”仍是待测假设；
原稿没有提供这些新算法的实测结果。
}}
\vfill
\noindent\textbf{交付范围}\quad 算法分析、逐步公式推导、资源估算、重构优先级与实验设计。\\
\textbf{实施范围}\quad 未修改项目源码、配置或机器人运行状态；仅生成本报告及其可编辑源文件。\\
\textbf{阅读建议}\quad 先读第 1、2、8、11 节获得决策框架，再按需阅读各算法推导。
\end{titlepage}
\tableofcontents
\clearpage
\section{结论、审阅对象与优先级}
\subsection{直接回答：哪些部分值得加速或重构}
\textbf{若目标是相同在线耗时下取得更好的导航效果，首选“GPU 批量控制评估 + 多场景不确定性模型 + 可复用几何预计算”，其次是“GPU 多假设定位与扫描匹配”。}
这两条路线分别增加控制问题覆盖的未来情况与定位问题保留的候选解，能改变算法能力上限。
单独加速现有直连碰撞只缩短一个布尔查询，收益取决于该查询在闭环中占据多少时间。

原稿的控制器并非每周期都运行 MPPI：直接末端收敛与几何路径跟踪优先，无法使用时才进入 MPPI。
因此 MPPI 虽有最清晰的规则并行结构，却不一定是整机最大耗时来源。
固定地图运行应优先分析 AMCL；在线建图才分析 Karto 与位姿图。
不能将不同模式的收益简单相加。

\begin{longtable}{p{0.18\textwidth}p{0.23\textwidth}p{0.31\textwidth}p{0.17\textwidth}}
\caption{针对本系统的建议。优先级是研究排序，不是实测性能排序。}\\
\toprule 对象 & 原始规模/结构 & 建议与性能收益机制 & 建议优先级\\\midrule\endfirsthead
\toprule 对象 & 原始规模/结构 & 建议与性能收益机制 & 建议优先级\\\midrule\endhead
MPPI & $1600$ 条、$56$ 步、七项代价 & 保留基线后移植批评估；再扩大为多模式、多场景控制，减少遗漏可行行为和低估风险 & 高；看分支占比\\
AMCL & $700$--$3000$ 粒子，至多 $120$ 束 & 并行粒子束评分；进一步维护多峰分布与多分辨率位姿网格 & 固定地图模式高\\
相关扫描匹配 & 候选位姿 $\times$ 激光点 & GPU 穷举或有上界的分支定界；多起点连续精化，提高搜索覆盖与亚栅格拟合能力 & 在线 SLAM 高\\
足迹碰撞 & 位姿 $\times$ 障碍 $\times$ 边 & 距离场安全筛选、空间索引、方向配置空间；把工作量从每次查询转到地图更新 & 共享基础设施高\\
制动扫掠 & 最多 $33$ 条候选速度轨迹 & 批量检查全部候选，或增加动态/执行误差情景；保持原接纳顺序与失败处理 & 中；需低尾延迟\\
Smac/航向修复 & 优先队列；每层至多 $4$ 航向 & GPU 做边验证/配置空间查表；扩大候选航向族，CPU 保留搜索管理 & 中\\
距离场/地图构造 & 局部图 $400\times400$ & EDT、膨胀和射线统计批处理；多个消费者复用，但语义独立 & 中到高；看复用率\\
位姿图优化 & 稀疏图，线程数 $2$ & 大图再考虑 GPU 稀疏求解；小图优先改善模型、异常边处理 & 条件性中低\\
EKF/四元数/仲裁 & $15$ 维矩阵及常数级逻辑 & 保留 CPU；GPU 不增加可观测性，单次规模难覆盖调度成本 & 低\\
\bottomrule
\end{longtable}

\subsection{输入、证据与实际核实范围}
审阅输入为用户提供的 \code{duojin02_algorithm_manuscript_source.zip}，含 11 个 \LaTeX{} 文件、构建脚本及验证/源码清单，共 16 个条目。
压缩包 SHA-256 为：
\begin{quote}\small\ttfamily
f7a9797307c3cca22b929a564263ace6dc42e86e7bff84caa2af9f5400c8e385
\end{quote}
本次读取全部正文、附录及原生代价说明，并逐项比对源清单：其中 101 个文件均存在，散列全部相同。
另只读检查当前 GPU 几何实现、直连调用点、制动扫掠与配置。
清单对应 Git 基点为 \code{89c2f3acae1d05e804b31f720fa52757b5ef190d}，但工作区有既存未提交修改，实际以散列为准。
主机查询到 RTX 5070、12227 MiB 显存、驱动 580.178.04；\textbf{这只是设备识别，未执行新 CUDA 性能实验}。

原稿 \code{evidence.tex} 报告 14472 个几何查询、470 个实际 CUDA 批次、439 个 CPU 边界回查，最终零差异。
这是既有归档记录，本报告未重新执行这些实验。
同文还记载较早产物中一个 $8000$ 障碍、$256$ 位姿的早停负载：
\[
T_C=0.09295\ \mathrm{ms},\qquad T_G=0.89514\ \mathrm{ms},\qquad
T_G/T_C=9.6303.
\]
这些数字说明“并行工作总量大”不充分，不构成现版本或新方案的速度预测。

本文把陈述分为三类：
\begin{itemize}
\item \textbf{已实现事实}：来自原稿及相同散列的源码；明确保留版本、数值与分支语义。
\item \textbf{数学结论}：在写出的条件下由公式推导；实数证明与浮点实现验证分开。
\item \textbf{设计建议}：尚未实现的新计算结构；规模数值是预算例，不是实测吞吐或精度。
\end{itemize}
外部方法只引用作者论文、官方文档或版本化官方源码；链接与访问日期见文末。
较新 Ceres 文档的 CUDA 能力不代表本项目 Humble 依赖已经支持。

\subsection{统一记号与量纲}
\begin{longtable}{p{0.24\textwidth}p{0.66\textwidth}}
\toprule 符号 & 定义\\\midrule\endfirsthead
\toprule 符号 & 定义\\\midrule\endhead
$q=(p,\theta)$ & 平面位姿，$p\in\R^2$，单位 m；$\theta$ 单位 rad\\
$R(\theta),J$ & $R=\left[\begin{smallmatrix}\cos\theta&-\sin\theta\\\sin\theta&\cos\theta\end{smallmatrix}\right]$，$J=\left[\begin{smallmatrix}0&-1\\1&0\end{smallmatrix}\right]$，$R'=RJ$\\
$\mathcal F,r_F,N_v$ & 车体坐标系内填充凸足迹、最大点半径、顶点数；$\mathcal F(q)=p+R(\theta)\mathcal F$\\
$\mathcal O,N_o,\rho_j$ & 障碍圆盘集合、障碍数、第 $j$ 个圆盘半径\\
$N_q,N_p,B$ & 查询位姿数、定位粒子数、有效激光束数；三者不混用\\
$K,H,M$ & 控制候选轨迹数、预测步数、不确定性场景数\\
$h,\delta,r_g$ & 积分步长（s）、碰撞余量（m）、栅格边长（m）\\
$N_x,N_y,N_\theta$ & 网格横纵单元数与方向层数；$G=N_xN_y$\\
$D_c,D_O$ & 到占据栅格中心的距离、到实际建模障碍集合的距离（m）\\
$S_n,\lambda,w_n$ & 第 $n$ 条控制轨迹总代价、同标度温度、无量纲归一化权重\\
$P,\Sigma$ & 协方差；须依状态/噪声单位解释，不能与旋转矩阵混用\\
$T_0,T_*,f,s$ & 基线端到端时延、目标时延预算、可优化时间比例、对应子任务加速比\\
$\epsilon$ 下标 & 各类误差界；数值误差、采样误差、建模误差必须各自标注\\
\bottomrule
\end{longtable}
本文的“代价”是具有明确权重标度的工程数值，不自动具有焦耳量纲。
位置与角度一起进入二次型时引入尺度长度 $\ell_\theta$：
$z=(p_x,p_y,\ell_\theta\theta)^\top$，则三个坐标均以 m 计。
对角度只在不跨主值分支的邻域求导。

\section{以固定在线时延为约束，而非追求孤立加速比}
\subsection{从端到端预算推导可用计算扩张量}
设基线周期总时延为 $T_0$，其中可重构部分为 $fT_0$，其余为 $(1-f)T_0$。
假设同工作量在 GPU 上的有效加速比为 $s$；新算法将这部分工作增加 $\alpha$ 倍。
令新增传输、同步、地图更新分摊及调度开销为 $T_{\rm oh}$。
则简化的串行关键路径模型为
\begin{equation}
T_{\rm new}=(1-f)T_0+\frac{\alpha fT_0}{s}+T_{\rm oh}.
\label{eq:budget}
\end{equation}
要求 $T_{\rm new}\le T_0$，先减去 $(1-f)T_0$，再除以 $fT_0/s>0$：
\begin{align}
\frac{\alpha fT_0}{s}+T_{\rm oh}&\le fT_0,\\
\alpha&\le s\left(1-\frac{T_{\rm oh}}{fT_0}\right).
\label{eq:expansion}
\end{align}
若 $T_{\rm oh}\ge fT_0$，即使计算无限快也无法在这个模型中获得正的扩张预算。
若仅搬迁而不扩大问题，即 $\alpha=1$ 且忽略新增开销，则
\[
\frac{T_0}{T_{\rm new}}=\frac1{1-f+f/s},
\]
这就是 Amdahl 型限制。

\textbf{纯示例：}若 $T_0=40$ ms、$f=0.75$、有效 $s=12$、$T_{\rm oh}=2$ ms，
则 $\alpha\le11.2$。把 GPU 部分扩大 $8$ 倍后，模型给出
$T_{\rm new}=10+8\times30/12+2=32$ ms。
这说明用户提出的换取方式在数学上成立；但 $s$ 必须是\emph{扩大后对应负载上的有效吞吐}。
缓存命中、寄存器溢出、显存访问与归约结构变化会使 $s$ 随 $\alpha$ 改变，不能线性外推到任意规模。

\subsection{将“算法更好”写成可检验优化问题}
设设计向量
$d=(K,H,M,N_p,N_\theta,r_g,\text{搜索层数},\text{模型复杂度})$，
质量指标向量 $Q(d)$ 包括定位错误率、到达率、尾部碰撞风险、跟踪误差与控制平滑度。
正确研究问题是
\begin{equation}
\max_d Q(d)\quad\text{受限于}\quad
T_{99}(d)\le T_*,\quad
\Pr(T(d)>T_*)\le\eta_T,\quad
V(d)\le V_{\max}.
\label{eq:pareto}
\end{equation}
多指标没有天然唯一全序，应报告 Pareto 前沿，或预先固定主要指标与次要约束。
这里 $T_{99}$ 为完整调用路径 P99，$V$ 是峰值显存；$\eta_T$ 是允许超期比例。
当前控制周期为 50 ms，但“比原 CPU 不慢”应同时满足原基线时延标准。
更快返回却使用更老地图，也不必然改善控制。

\subsection{数据驻留、更新频率与真正摊销}
对 GPU 单任务，分解
\begin{equation}
T_G=T_{\rm prepare}+T_{\rm H2D}+T_{\rm launch}
+T_{\rm kernel}+T_{\rm D2H}+T_{\rm wait}+T_{\rm recheck}.
\end{equation}
设每幅地图建立索引耗时 $T_{\rm build}$，两次更新之间有 $Q_{\rm reuse}$ 次查询，
每查询核耗时 $t_G$，则平均成本是
\begin{equation}
\bar T_G=\frac{T_{\rm build}}{Q_{\rm reuse}}+t_G.
\label{eq:amort}
\end{equation}
若 CPU 单查询成本为 $t_C$，需满足
$Q_{\rm reuse}>T_{\rm build}/(t_C-t_G)$，且分母为正。
这不意味着地图更新峰值也被平均化：关键周期仍可能承受完整 $T_{\rm build}$。
双缓冲可隐藏更新，但必须带地图版本、时间戳与最大允许年龄；不能默默使用过期障碍图。

局部代价图 $4\mathrm m\times4\mathrm m$、分辨率 $0.01$ m，包含 $160000$ 格；
一个 float 距离场约 $0.64$ MB（十进制 MB）。
静态 AMCL 距离场可长驻，动态局部图不能假设长期不变。
最终保护的激光记忆、局部 costmap 与全局定位地图具有不同未知区、时间戳和障碍语义；
可以共享 GPU 服务和存储模式，不能直接合并为同一个无标签数组。
这与 NVIDIA 对最小化传输及保持设备数据驻留的建议一致\cite{cuda}。

\subsection{分支调用率与过期输入的质量代价}
设每周期调用 MPPI 的比例为 $\pi_{\rm M}$，
其他路径耗时 $T_{\rm D}$，MPPI 路径耗时 $T_{\rm M}$，则平均时间
\[
\E T=(1-\pi_{\rm M})\E T_{\rm D}+\pi_{\rm M}\E T_{\rm M}.
\]
降低 $T_{\rm M}$ 的平均收益被 $\pi_{\rm M}$ 缩放；P99 仍可能由 MPPI 决定。
如果新增算法让 MPPI 在以前直接跟踪的周期也运行，必须把新调用率计入预算。
可以只在狭窄、歧义或高不确定性时启用重型算法，而不能用旧分支占比估算新系统。

若观测到控制输出的总年龄为 $\tau$，短时状态误差满足
$\|\Delta p\|\lesssim v\tau$、$|\Delta\theta|\lesssim|\omega|\tau$，
则足迹点误差保守近似为
\[
\epsilon_{\rm age}\lesssim(v+r_F|\omega|)\tau.
\]
因此增加计算量必须维持时效。吞吐更高但排队更久，可能抵消更丰富模型的收益。

\section{几何重构：从反复遍历障碍到可复用查询结构}
\base{geometry.tex；control.tex 的制动扫掠；gpu/geometry.cu}
\subsection{现有算法的运算对象与早停优势}
对第 $i$ 个位姿与第 $j$ 个障碍，把中心变换到车体坐标：
\[
y_{ij}=R(\theta_i)^\top(o_j-p_i).
\]
第 $v$ 条边为 $e_v=a_{v+1}-a_v$，点到线段距离为
\begin{align}
t_v&=\clip\left(\frac{(y_{ij}-a_v)^\top e_v}{e_v^\top e_v},0,1\right),\\
d_v&=\norm{y_{ij}-a_v-t_ve_v}.
\end{align}
在包围盒门控 $A_{ij}$ 内，局部碰撞是
\[
c_{ij}=A_{ij}\land
\left\{\left[\bigwedge_v e_v\times(y_{ij}-a_v)\ge-\epsilon_\times\right]
\lor\left[\min_v d_v\le\delta+\rho_j\right]\right\}.
\]
总结果 $C_i=\bigvee_jc_{ij}$。理想算术下，OR 的交换律、结合律允许任意线程划分；
这只改变求值顺序，未改变算法复杂度 $O(N_qN_oN_v)$。

CPU 对某个位姿首次遇到碰撞就停止。
用简化独立模型表示每次障碍检测碰撞概率为 $p_{\rm hit}$，
令 $L$ 为至多检查 $N_o$ 个障碍时的实际检查次数。
利用正整数随机变量的尾和公式：
\begin{align}
\E L&=\sum_{\ell=1}^{N_o}\Prob(L\ge\ell)
=\sum_{\ell=1}^{N_o}(1-p_{\rm hit})^{\ell-1}\\
&=\frac{1-(1-p_{\rm hit})^{N_o}}{p_{\rm hit}},\quad p_{\rm hit}>0.
\label{eq:early}
\end{align}
若 $p_{\rm hit}\to0$，极限为 $N_o$；若较大，趋近 $1/p_{\rm hit}$。
实际障碍并不独立，此式只解释机制。
此外直连 CPU 路径一旦某个位姿碰撞，会停止整条路径；GPU 当前把整批位姿都算完。
因此仅按 $N_qN_o\ge262144$ 分流忽略了两级早停。
当前实现已在主机预计算位姿正余弦，并按场景版本复用上传；
这些不能再作为“尚未实现”的新加速点。

\subsection{先缩小问题：空间桶、排序与候选合批}
对每个位姿，只需要考虑与扩张足迹包围盒相交的障碍。
若空间索引返回 $n_i$ 个邻近障碍，则窄相位工作量变为
\[
O\!\left(N_q\,c_{\rm index}+N_v\sum_i n_i\right),
\]
其中 $c_{\rm index}$ 是查询索引成本，均匀桶在局部密度有界时可能近似常数；
聚集退化情况下仍会接近全遍历。桶边长是算法参数，应按足迹范围和障碍密度选取。
CPU 也应使用同样的候选削减作为公平基线，而不能只比较 GPU 优化算法和 CPU 暴力算法。

多个横向居中候选可把全部直连位姿合成一次 GPU 调用。
但原控制器先按 $\alpha_c\in\{1,0.65,0.4\}$ 分组，再按净空分数稳定排序，
选择第一个通过候选。正确合批步骤是“并行计算全部可行布尔值，再按原顺序找首个真值”，
不是在 21 个候选里重新选全局最优。
重构成全局最优属于另一算法，需单独比较。

\subsection{距离场安全筛选：推导可保证的下界}
定义障碍集合 $O$ 与其非负距离场
$D_O(x)=\inf_{o\in O}\|x-o\|$。
对任何 $x,y$，三角不等式给出
\[
\|x-o\|\le\|x-y\|+\|y-o\|.
\]
对 $o$ 取下确界得 $D_O(x)\le\|x-y\|+D_O(y)$；
交换 $x,y$，于是
\begin{equation}
|D_O(x)-D_O(y)|\le\|x-y\|.
\label{eq:lipschitz}
\end{equation}
任意足迹点 $a\in\mathcal F$ 满足 $\|R(\theta)a\|\le r_F$。
在式 \eqref{eq:lipschitz} 取 $x=p+R(\theta)a,y=p$：
\[
D_O(p+R(\theta)a)\ge D_O(p)-r_F.
\]
再对所有 $a$ 取最小，得到足迹到障碍的距离下界
\begin{equation}
\dist(\mathcal F(q),O)\ge D_O(p)-r_F.
\label{eq:centerbound}
\end{equation}
所以 $D_O(p)>r_F+\delta$ 是充分安全条件；不满足只表示“不确定”，不能直接判碰撞。
中心圆筛选在宽阔区域很有效，在窄道会保守，应落回完整填充足迹检查。

实际 EDT 常给到占据格\emph{中心}的距离 $D_c$。
每个边长 $r_g$ 的占据格均位于以中心为圆心、半径 $\rho_g=\sqrt2r_g/2$ 的圆盘内。
对任意格内点 $o$、格中心 $c$，
$\|x-o\|\ge\|x-c\|-\rho_g$，
所以
\begin{equation}
D_O(x)\ge\max(0,D_c(x)-\rho_g).
\end{equation}
若在最近查询格中心 $x_g$ 查表，$\|x-x_g\|\le\rho_g$；
表值的绝对数值误差上界为 $\epsilon_D$，
则
\begin{equation}
\underline D_O(x)=
\max\{0,\widehat D_c(x_g)-\epsilon_D-2\rho_g\}\le D_O(x).
\label{eq:edtbound}
\end{equation}
第一个 $\rho_g$ 是障碍格几何外扩，第二个是查询量化误差，不能重复遗漏或混为一个。
若用四格双线性插值，权重 $a_s\ge0,\sum_s a_s=1$，由 Lipschitz 性
\[
\left|\sum_s a_sD_c(x_s)-D_c(x)\right|
\le\sum_s a_s\|x_s-x\|\le\sqrt2r_g;
\]
可以使用这个较松但明确的插值误差界，再加 $\rho_g$ 与数值误差。
纹理插值值本身不自动是安全下界。

\textbf{与源码容差的区别。}
上述保证针对真实填充多边形及建模障碍。
CPU 内点判定使用 $\epsilon_\times=10^{-10}\,\mathrm m^2$，可能把近边界外点也判入。
若要求连同此容差的布尔输出也一致，可定义
\[
\mathcal F_\epsilon=\{y:e_v\times(y-a_v)\ge-\epsilon_\times,\ \forall v\}
\]
并使用其外接半径 $r_{F,\epsilon}\ge r_F$ 做筛选。
有效非退化凸足迹的该集合有界，去除冗余约束后，其顶点由相邻平移支持线求交得到。
于是把式 \eqref{eq:centerbound} 中 $r_F$ 换为 $r_{F,\epsilon}$，
再计入距离阈值，便不忽略容差假内点。
浮点上下界还必须向保守方向舍入；无经验证误差界的点继续回查 CPU。

\subsection{多圆覆盖使距离场适应非圆车体}
设离线构造有限圆盘，使
$\mathcal F\subseteq\bigcup_{a=1}^{A}\mathbb B(c_a,r_a)$。
第 $a$ 个圆盘变换后圆心为 $p+R(\theta)c_a$。
若每个覆盖圆都满足
\begin{equation}
\underline D_O(p+R(\theta)c_a)>r_a+\delta,\qquad a=1,\ldots,A,
\label{eq:cover}
\end{equation}
则覆盖并集与障碍隔离，子集 $\mathcal F(q)$ 也隔离。
每个位姿只需 $A$ 次距离场查表；$A$ 可以远小于 $N_oN_v$。
圆盘越多越贴合足迹，窄道保守性越小，但查询增加。
该方法是“显存与预处理换在线工作”的具体例子。

只检查足迹顶点并不成立：小障碍可以完全处在足迹内部，距所有顶点都很远。
只检查边界采样也不能自动解决内部障碍问题。
必须使用完整覆盖、精确多边形，或后述配置空间占据。
对变半径激光圆盘，不能直接套用等半径占据格 EDT；
可先保守栅格化所有圆盘并计入栅格化误差，或用保留半径的空间索引精查。

\subsection{精确 EDT 的可分离结构，而非近似距离的盲用}
令二值图的代价函数 $f(i,j)$ 在障碍中心取 $0$，其他格取 $+\infty$。
以格为单位的平方 EDT 为
\[
E(x,y)=\min_{i,j}\{(x-i)^2+(y-j)^2+f(i,j)\}.
\]
有限集合最小化可交换次序：
\begin{align}
g(i,y)&=\min_j\{(y-j)^2+f(i,j)\},\\
E(x,y)&=\min_i\{(x-i)^2+g(i,y)\},\qquad D_c=r_g\sqrt E.
\end{align}
一维变换是抛物线 $\varphi_i(x)=(x-i)^2+g_i$ 的下包络。
两条 $i\ne j$ 抛物线交点由
$x^2-2ix+i^2+g_i=x^2-2jx+j^2+g_j$ 得
\[
x_{ij}=\frac{j^2+g_j-i^2-g_i}{2(j-i)}.
\]
维护下包络交点的有序栈，每条抛物线至多进栈、出栈一次，
可得一维线性总工作；二维逐轴处理为 $O(G)$。
行列独立，但一行内的包络维护不是全部元素无依赖。
PBA/PBA+ 通过分带及合并提供 GPU 精确 EDT 路线\cite{edt,pba}。
这里“精确”指选定离散中心距离，不是物理环境距离。

普通 Jump Flooding 一类近似方法若漏掉最近障碍，返回较远种子距离，即\emph{高估}间隙；
高估不能直接用于安全放行。只有取得可证明误差上界并按式 \eqref{eq:edtbound} 扣除，
或把它仅用于排序再精查，才符合本报告的保证条件。

\subsection{更深入的替换：位置--航向配置空间查表}
定义 Minkowski 和 $A\oplus B=\{a+b:a\in A,b\in B\}$。
在固定方向 $\theta$ 下，足迹与障碍相交等价于
\begin{align}
(p+R(\theta)\mathcal F)\cap O\ne\varnothing
&\iff \exists f\in\mathcal F,o\in O:\ p+R(\theta)f=o\\
&\iff p=o-R(\theta)f\\
&\iff p\in O\oplus[-R(\theta)\mathcal F].
\end{align}
因此可预计算 $N_\theta$ 个配置空间占据切片：
\begin{equation}
C_\theta(p)=\ind\{p\in O\oplus[-R(\theta)\mathcal F]\oplus\mathbb B(0,\delta)\}.
\label{eq:cspace}
\end{equation}
把复杂几何判断变成三维表的查询，尤其适合大量轨迹与方向候选共用一幅地图。

若方向均匀采样，间隔 $\Delta\theta=2\pi/N_\theta$，最近方向误差至多 $\Delta\theta/2$。
对足迹点，
\begin{align}
\|R(\theta)a-R(\theta_j)a\|
&=2\|a\|\left|\sin\frac{\theta-\theta_j}{2}\right|\\
&\le2r_F\sin\frac{\Delta\theta}{4}
=:e_\theta.
\label{eq:anglebound}
\end{align}
查询中心量化误差为 $e_p\le\sqrt2r_g/2$。
把建表余量增为 $\delta+e_p+e_\theta$，
再加障碍栅格化/数值误差，最近层的“无碰撞”可成为真实查询的充分安全条件。
否则直接取最近方向层可能漏碰撞。
若非等间隔方向，应使用当地相邻间隔或最大间隔，不能假定原 Smac 的 16 个方向均匀。

取说明性 $r_F=0.30$ m、$N_\theta=64$，则 $e_\theta\approx14.72$ mm；
$1$ cm 网格还有 $e_p\approx7.07$ mm。
这远大于原直连余量与采样共同形成的 8 mm 名义下界。
所以 64 层保守查表可能封死窄道，合理方案是“宽区直接通过，近边界更细方向或精确检查”，
并非粗查表完全替代原几何。

对 $400\times400$ 图、64 方向：状态总数 $10.24$ 百万；
每状态 1 bit、1 byte、float 分别约 $1.28$、$10.24$、$40.96$ MB。
双缓冲至少翻倍。若足迹模板含 $A_F$ 个格，直接膨胀建表工作为
$O(N_\theta GA_F)$；内存很小不意味着构建很快。
可用并行形态学、位运算或离线静态地图预计算。
动态障碍更新只影响其与反射足迹模板的作用区；障碍删除需要计数或局部重建，
不能只 OR 新障碍，否则旧障碍永远残留。
滚动窗口平移时还须维护世界坐标原点和新进入区域，不能把数组循环移动当成完整更新。

\subsection{批量制动与连续间隙：能增加什么检查}
原防护先检查实测速度，再按 8 个缩放值，每个检查请求及 3 个过渡速度；
最坏有 $1+8\times4=33$ 条制动轨迹。
可先独立生成每条轨迹，再把其采样位姿合批查询，并按轨迹 OR 归约。
令 $b_{s,\beta}$ 是轨迹碰撞标志，
\[
s_*=\max\{s:\bigvee_{\beta\in\{1,0.25,0.5,0.75\}}b_{s,\beta}=0\}.
\]
这能保留原“按比例选最大可接受值”语义，且无需假定可行性对 $s$ 单调。
不同轨迹的 $h_b$、终止步数不同，批内需偏移数组；
原算法先推进再减速的顺序不能改成连续制动距离公式。

更进一步，可对执行迟滞、减速度、初始位姿及障碍速度生成多场景，
检查量从 33 条扩到 $33M$ 条；这是额外计算换更充分模型的入口。
但每条情景仍是模型假设，其覆盖性见第 \ref{sec:deep} 节。

对静态障碍，原直连长度参数
$L=\|\Delta p\|+r_F|\Delta\theta|$，
取 $N=\lceil L/\ell_s\rceil$。
任意插值点距最近样本参数差至多 $1/(2N)$，
足迹 Hausdorff 距离至多 $L/(2N)\le\ell_s/2$。
再由集合距离 Lipschitz 性，
\begin{equation}
d_{\rm continuous}>\delta-\ell_s/2.
\end{equation}
代入 $\delta=10$ mm、$\ell_s=4$ mm 得 8 mm；
增加采样会提高该\emph{理想模型}下界。
若实际位姿、障碍和时延误差界分别为
$\epsilon_p+r_F\epsilon_\theta$、$\epsilon_o$、$v_{\rm rel}\tau$，
则保证余量应再扣除它们。
GPU 不能靠更密几何采样消除未建模的定位或执行误差。

\section{MPPI：从并行采样到更充分的控制优化}
\base{control.tex 的 MPPI 部分；native_costs.tex}
\subsection{先明确基线究竟在算什么}
当前 $K=1600,H=56,h=0.05$ s，预测跨度 $Hh=2.8$ s，
每次优化包含 $KH=89600$ 个预测状态、$3KH=268800$ 个控制分量。
float 控制数组约 1.0752 MB，float 三维位姿数组同样约 1.0752 MB；
噪声、速度、路径索引和临时结果需另计。

名义控制为 $\bar u_\ell=(\bar v_{x,\ell},\bar v_{y,\ell},\bar\omega_\ell)$，
扰动 $\epsilon_{n,\ell}\sim\mathcal N(0,\Sigma_u)$，
$\Sigma_u=\diag(0.15^2,0.12^2,0.18^2)$。
实际采样控制 $\tilde u_{n,\ell}=\bar u_\ell+\epsilon_{n,\ell}$。
原配置复用噪声，不能把每周期重新生成独立噪声描述为等价移植。
Omni 模型第一步用实测速度，其后用前一列采样控制：
\begin{equation}
\hat u_{n,0}=u^{\rm meas},\qquad
\hat u_{n,\ell}=\tilde u_{n,\ell-1}\quad(\ell\ge1).
\end{equation}
随后
\[
p_{n,\ell+1}=p_{n,\ell}+hR(\theta_{n,\ell})\hat v_{n,\ell},\quad
\theta_{n,\ell+1}=\theta_{n,\ell}+h\hat\omega_{n,\ell}.
\]
第一预测状态因而对所有轨迹相同。输出序列还有约束、滤波和移位，不能忽略。

总代价
\begin{equation}
S_n=\sum_r\mathcal J_r+
\gamma\sum_{\ell,d}\frac{\bar u_{\ell,d}\epsilon_{n,\ell,d}}{\sigma_d^2},
\qquad \gamma=0.015.
\label{eq:mppicost}
\end{equation}
式中 $r$ 遍历启用的代价评估器，$\mathcal J_r$ 是其返回的标量代价。七个启用 critic 含路径索引、阈值和碰撞分支，不是统一二次型。
当前原生障碍 critic 的轮廓检查与项目填充几何不同；
将其替换成填充碰撞是模型增强，不能计为纯 CUDA 后端改写\cite{mppisrc}。

\subsection{指数权重为何是合理的，完整推导}
对固定采样集考虑
\[
\min_{w\ge0,\ \sum_nw_n=1}
F(w)=\sum_nw_nS_n+\lambda\sum_nw_n\log(Kw_n),\qquad\lambda>0.
\]
第二项是相对均匀分布的熵正则。
拉格朗日函数
$\mathcal L=F(w)+\mu(\sum_nw_n-1)$。
在正权重区域逐项求导：
\[
\frac{\partial\mathcal L}{\partial w_n}
=S_n+\lambda[\log(Kw_n)+1]+\mu=0.
\]
移项、指数化得
\[
w_n=\frac1K\exp[-S_n/\lambda]\exp[-1-\mu/\lambda].
\]
将其代入 $\sum_nw_n=1$ 消去共同常数：
\begin{equation}
w_n=\frac{\exp[-(S_n-S_{\min})/\lambda]}
{\sum_j\exp[-(S_j-S_{\min})/\lambda]}.
\label{eq:softmax}
\end{equation}
Hessian 是 $\diag(\lambda/w_n)$，对内部权重严格正定，
因此驻点是该有限单纯形问题的唯一最优解。
其中 $S_{\min}=\min_n S_n$；减去 $S_{\min}$ 只抵消共同乘数，不改变实数结果。

令 $\mathcal U$ 为当前配置及动态限速共同确定的三轴速度盒，$\clip_{\mathcal U}$ 为逐轴裁剪。更新
\begin{align}
\bar u_\ell^+&=\clip_{\mathcal U}\left(\sum_nw_n\tilde u_{n,\ell}\right)\\
&=\clip_{\mathcal U}\left(\bar u_\ell\sum_nw_n+\sum_nw_n\epsilon_{n,\ell}\right)\\
&=\clip_{\mathcal U}\left(\bar u_\ell+\sum_nw_n\epsilon_{n,\ell}\right).
\end{align}
这只证明固定样本的权重最优，不证明整个非线性轨迹空间全局最优。
路径积分/信息论控制的总体背景见作者论文\cite{itmppi}；
本文不把该背景直接当成原生各 critic 的理论最优性证明。

\subsection{并行结构：轨迹独立，部分时间递推还可转成扫描}
保守第一版可令每个线程或线程组负责一条轨迹，顺序完成 $H$ 步。
有 $K$ 份独立工作，之后做最小值、指数求和与加权控制归约。
也可分开轨迹生成与状态评分，提高地图读取并行度。

对当前开环 Omni 欧拉模型，时间方向并非完全不能并行。
由于 $\hat\omega_{n,\ell}$ 不依赖未来位置，
\begin{equation}
\theta_{n,\ell}=\theta_0+h\sum_{j=0}^{\ell-1}\hat\omega_{n,j}.
\end{equation}
先对角速度作 exclusive prefix sum，再独立计算
$d_{n,\ell}=hR(\theta_{n,\ell})\hat v_{n,\ell}$，最后
\begin{equation}
p_{n,\ell}=p_0+\sum_{j=0}^{\ell-1}d_{n,j}.
\end{equation}
于是两次并行前缀和即可产生全部轨迹。
工作仍是 $O(KH)$，理想并行深度从 $O(H)$ 降到 $O(\log H)$。
但 $H=56$ 很短，扫描的同步与临时数据可能比顺序寄存器积分更贵。
采用依赖状态的动力学、反馈策略或接触模型后，此拆解一般不成立。
并行求和也改变舍入顺序，须与确定噪声的 CPU 基线逐步比较。

路径 critic 存在整批依赖，例如先求所有轨迹终点的最远匹配路径索引，
其他 critic 才知道目标索引，因而需要分阶段执行，不能把七项全部随意融合。
末尾 Savitzky--Golay 滤波目前按原地更新顺序执行，
它不是可直接替换的普通并行卷积。
路径关联存在原稿指出的未定义末端分支，移植前应明确合法输入域与期望修复语义；
GPU 不能替一个未定义行为建立等价证明。

\subsection{扩大采样为何可能有效，又为何不足够}
设随机生成一条候选落入“质量合格集合” $\mathcal A$ 的概率为 $p_A$，
不同候选独立，则全部 $K$ 条未命中概率为
\[
\Prob(\forall n,\ U_n\notin\mathcal A)=(1-p_A)^K.
\]
因而至少一次命中概率
\begin{equation}
P_{\rm cover}(K)=1-(1-p_A)^K.
\label{eq:coverage}
\end{equation}
若 $p_A=10^{-3}$，$K=1600$ 时约 $79.83\%$；
$K=6400$ 时约 $99.83\%$。
这是指定采样分布下的覆盖计算，不是导航成功率。
且对嵌套样本集有
$\min_{n\le K'}S_n\le\min_{n\le K}S_n$（$K'\ge K$），
但 MPPI 输出是加权平均，并非最优样本本身，控制效果不必随 $K$ 单调。

权重有效样本数
\[
K_{\rm eff}=\frac1{\sum_nw_n^2}.
\]
均匀权重时 $K_{\rm eff}=K$，一个权重占优时接近 $1$。
更大 $K$ 若仍只集中在同一错误模式，不能解决多模态绕障。
高维独立噪声很难产生持续一致的侧向绕行，增加结构化候选可能比无差别扩样更有效。

\subsection{结构化控制序列：把预算花在更有意义的方向}
原控制自由度为 $3H=168$。
令 $B\in\R^{H\times L_c}$ 为固定时间基函数矩阵，
例如分段线性或样条；$L_c<H$。
系数 $C\in\R^{L_c\times3}$ 决定控制
\begin{equation}
U=BC,\qquad u_\ell=\sum_{a=1}^{L_c}B_{\ell a}C_a.
\label{eq:basis}
\end{equation}
对系数采样 $\Delta C$，即得到时间相关控制扰动。
若系数协方差为 $\Sigma_C$，
按列向量化后控制协方差为
\[
\Cov(\operatorname{vec}\Delta U)=
(I_3\otimes B)\Sigma_C(I_3\otimes B)^\top.
\]
它一般低秩，因此不能把原满秩独立控制噪声的逆协方差公式原样套过来。
合理做法是在系数空间定义采样与熵正则更新，或重新推导提议分布重要性修正。

这使每个样本更容易形成平滑、连续的绕障动作，扩大相同样本数对“有意义控制”的覆盖。
代价是只能表达 $\operatorname{col}(B)$ 内的轨迹，可能遗漏急转向。
可以使用低频基函数为主、局部高频扰动为辅，并按当前状态增加停车、左绕、右绕等行为模式。
若改变采样分布 $q(U)$，而目标仍是
$p^*(U)\propto\exp[-J(U)/\lambda]p_0(U)$，
正确样本权重为
\begin{equation}
w_n\propto\exp[-J(U_n)/\lambda]\frac{p_0(U_n)}{q(U_n)},
\quad U_n\sim q,\quad q>0\text{ 于目标支持集}.
\label{eq:is}
\end{equation}
该式由重要性恒等式
$\int f p^*=\int f(p^*/q)q$ 得到。
混合行为采样须使用实际混合密度，不是随意给新样本沿用旧权重。
低秩基函数若改变支持集，实质上是受限控制问题，需按新的测度定义目标，不能假装保持全空间等价。

\subsection{更好的运动预测，比单纯增加随机样本更有针对性}
对于一步内恒定机体速度 $v$ 和角速度 $\omega$，
精确积分为
\begin{align}
p_{\ell+1}-p_\ell
&=R(\theta_\ell)\int_0^h R(\omega t)v\,dt\\
&=R(\theta_\ell)
\left[\frac{\sin(\omega h)}{\omega}I+
\frac{1-\cos(\omega h)}{\omega}J\right]v.
\label{eq:se2exact}
\end{align}
因为 $R(\omega t)=\cos(\omega t)I+\sin(\omega t)J$，逐项积分即可得到上式。
令 $a=\omega h/2$，还可写成数值稳定形式
\[
p_{\ell+1}-p_\ell=h\,\operatorname{sinc}(a)R(\theta_\ell+a)v,
\qquad \operatorname{sinc}(a)=\sin a/a,\quad\operatorname{sinc}(0)=1.
\]
原欧拉项为 $hR(\theta_\ell)v$。
由 $\sin x=x-x^3/6+\cdots$、$1-\cos x=x^2/2+\cdots$，
两者差的首项是
\[
\Delta p_{\rm exact}-\Delta p_{\rm Euler}
=\tfrac12\omega h^2R(\theta_\ell)Jv+O(h^3).
\]
也可不依赖级数，利用 $\|R(\omega t)-I\|\le|\omega|t$，
积分得单步界
\[
\|\Delta p_{\rm exact}-\Delta p_{\rm Euler}\|
\le\tfrac12|\omega|h^2\|v\|.
\]
说明性 $\|v\|=0.5$ m/s、$|\omega|=0.65$ rad/s、$h=0.05$ s 对应至多约 0.406 mm/步。
固定同一开环控制并保持角积分相同时，将 56 步三角不等式相加给出约 22.75 mm 的松上界，
不是实际预测误差测量。
这类替换无需巨大计算量，却可能比盲目增加 $K$ 更直接减少离散预测偏差。

进一步纳入一阶执行响应
$\dot v=(v_{\rm cmd}-v)/\tau_v$，
恒定输入一步的精确解为
$v^+=e^{-h/\tau_v}v+(1-e^{-h/\tau_v})v_{\rm cmd}$；
可对不同 $\tau_v$、侧滑系数和制动能力并行情景展开。
其参数必须用实测标定；未知动力学不能靠计算量凭空恢复。
固定 2.8 s 跨度而将步长减半，需将 $H$ 加倍；
只增加 $H$ 而保持 $h$ 不变，是延长未来而不是提高时间分辨率。

\subsection{代价误差如何放大到控制误差}
若各轨迹计算误差满足 $|\widehat S_n-S_n|\le\epsilon_S$，
则未归一化权重比例在 $[e^{-\epsilon_S/\lambda},e^{\epsilon_S/\lambda}]$ 内。
归一化分母也满足同样比例，因此
\[
e^{-2\epsilon_S/\lambda}\le\frac{\widehat w_n}{w_n}
\le e^{2\epsilon_S/\lambda}.
\]
逐项相减、取绝对值并求和：
\begin{equation}
\|\widehat w-w\|_1\le
\min\{2,e^{2\epsilon_S/\lambda}-1\}.
\label{eq:softbound}
\end{equation}
若采样控制满足 $\|\tilde u_{n,\ell}\|\le U_{\max}$，
则加权平均误差不超过 $U_{\max}\|\widehat w-w\|_1$；
盒裁剪在欧氏范数下不扩大这个差异。
高斯采样本身无全局有界 $U_{\max}$，该界应使用实际有限批次的最大范数，
或明确限制后的采样域。
当前 $\lambda=0.3$，不可把未经分析的低精度代价误差当成无关紧要。
地图索引取错一个栅格更会产生离散分支变化，不属于微小舍入扰动。

\section{定位：并行粒子评分与多假设保持}
\base{localization.tex 的 AMCL、地图一致性部分}
\subsection{现有 AMCL 的核心批量计算}
粒子 $q_i=(p_i,\theta_i)$，当前束 $b$ 的激光坐标点为
$s_b^L=r_b(\cos\varphi_b,\sin\varphi_b)^\top$。
令固定激光到车体外参为 $T^B_L=(R_{BL},t_{BL})$，则
$s_b^B=R_{BL}s_b^L+t_{BL}$，其中 $R_{BL},t_{BL}$ 明确将激光系点变换到车体系。
地图端点为
\[
e_{ib}=p_i+R(\theta_i)s_b^B.
\]
对截断中心距离场 $D_{\mathcal M}$，
原模型使用\cite{amclsrc}
\begin{equation}
\pi_{ib}=a\exp[-D_{\mathcal M}(e_{ib})^2/(2\sigma^2)]+b_0,
\quad a=0.85,\ \sigma=0.025\ {\rm m},\
b_0=0.15/r_{\max}.
\end{equation}
这里 $\pi$ 是原库保留工程标度的束评分，不重新补上高斯归一化系数。
粒子后验权重正比于
\begin{equation}
\widetilde w_i=w_i^{\rm prior}\prod_{b=1}^{B}\pi_{ib}.
\label{eq:amclprod}
\end{equation}
当前 $B\le120$、$N_p=700$--$3000$，
因此一次扫描最多 $84000$--$360000$ 个粒子束组合。
各组合共享地图、扫描与外参，是很适合 GPU 的二维批量运算。
对每个粒子算一次三角函数，把束点预先变换到车体系，再分别完成端点及查表。

在保持束筛选、越界策略、量化索引和评分完全一致时，
并行化不改变实数公式。
用纹理双线性读取替代原离散取格、增加束跳过或改变噪声宽度，
均属于观测模型变更，必须单独报告。
AMCL、MPPI 与碰撞可共享设备常驻的相关地图缓存，但必须分别保留其语义标签。

\subsection{稳定对数归一化与不能忽略的语义改变}
若 $w_i^{\rm prior}>0$，取对数：
\[
\ell_i=\log w_i^{\rm prior}+\sum_b\log\pi_{ib}.
\]
令 $m=\max_i\ell_i$，则
\begin{align}
w_i^+
&=\frac{e^{\ell_i}}{\sum_je^{\ell_j}}
=\frac{e^m e^{\ell_i-m}}{e^m\sum_je^{\ell_j-m}}\\
&=\frac{e^{\ell_i-m}}{\sum_je^{\ell_j-m}}.
\label{eq:logsum}
\end{align}
至少一个指数为 $1$，其余不大于 $1$，避免共同下溢。
先验为零的粒子记 $\ell_i=-\infty$；若全部非法或负无穷，必须定义失败处理。
在实数运算中等价，但原实现的权重全零时会回到均匀权重；
稳定形式可能避免该退化，所以\textbf{数值上可能改变旧输出且改善可靠性}，应标为数值算法升级。

不能将 $\sum_b\log\pi_{ib}$ 除以 $B$ 作为“稳定化”。
除法改变相对权重：
两粒子的似然比从 $\prod_b(\pi_{ib}/\pi_{jb})$ 变成其 $1/B$ 次方，
等于改变观测温度与统计模型，而不只是防溢出。
类似地，统一扩大 $\sigma$ 会改变地图匹配辨识度。

\subsection{重采样的并行边界与 KLD 语义}
设归一化权重为 $w_1,\dots,w_{N_p}$，累计分布
$F_j=\sum_{i=1}^jw_i$。
GPU prefix sum 可并行求出 $F$。
给定抽样数 $N'$，systematic 重采样取
$u_0\sim U[0,1/N')$、$u_m=u_0+m/N'$，
返回 $j(m)=\min\{j:F_j\ge u_m\}$。
二分查询可并行，工作约 $O(N'\log N_p)$，或利用有序合并做线性总工作。

但当前 KLD 重采样根据\emph{已经抽到的粒子}占据箱数 $\kappa$ 决定停止位置。
其停止长度与粒子序列相互依赖。
先算整批新粒子再按固定 $N'$ 取样，或者直接换 systematic，
会改变联合分布/停止规律，不是简单的并行求和等价替换。
初期可把运动采样与束评分搬到 GPU，保留 KLD、聚类与输出在 CPU；
若通信成为瓶颈，再单独设计分批 KLD 并声明与基线的差别。
原稿还指出 \code{pf_z=0.99} 被直接作为系数使用，不能据此宣称 99\% 置信保证。

\subsection{更大粒子数可减少什么误差}
为了给“更多计算换更好估计”一个明确条件，
令粒子 $X_i$ 独立来自提议分布 $q$，
未归一化重要性权重 $a_i=p(X_i)/q(X_i)$，
目标统计量为
$\mu=\int\phi(x)p(x)dx/\int p(x)dx$。
记
$A_N=N^{-1}\sum_i a_i\phi(X_i)$、
$B_N=N^{-1}\sum_i a_i$、
$B=\E_q a_i>0$。
自归一化估计为 $\hat\mu_N=A_N/B_N$。
在 $(A,B)$ 处一阶展开：
\begin{align}
\hat\mu_N-\mu
&\approx\frac{A_N-A}{B}
-\frac{A}{B^2}(B_N-B)\\
&=\frac1{NB}\sum_i a_i[\phi(X_i)-\mu],
\end{align}
所以在相应二阶矩有限、独立抽样及渐近条件下
\begin{equation}
\Var(\hat\mu_N)\approx
\frac{\Var_q(a[\phi(X)-\mu])}{NB^2}.
\label{eq:snis}
\end{equation}
粒子数增为 4 倍，标准差近似减半；模型偏差并不会减半。
经过重采样的粒子相关，不能机械按独立 $N$ 套用。
ESS 是退化诊断而非通用精度定理。

若某个真实可行模式在当前后验中质量为 $m_A$，
独立重采样 $N'$ 次一个也没保留的概率为
\[
P_{\rm lose}=(1-m_A)^{N'}.
\]
这解释扩大粒子预算有助于保留低质量但正确的候选。
但是如果初始提议根本不覆盖该模式，则 $m_A=0$，任意复制旧粒子都无济于事。
因此更深重构应同时保留全局探索或独立多模式分支，而非只增大一个收敛错误簇。

\subsection{更多激光束为何也有收益上限}
局部近似令残差
$r_b(q)=D_{\mathcal M}(p+R(\theta)s_b^B)$，
在场可微且小扰动下
\begin{equation}
j_b=\frac{\partial r_b}{\partial q}
=\nabla D_{\mathcal M}(e_b)^\top
\begin{bmatrix}I_2&R(\theta)Js_b^B\end{bmatrix}.
\label{eq:locjac}
\end{equation}
若残差是独立、方差 $\sigma_b^2$ 的近似高斯，
局部信息矩阵为
\[
\mathcal I=\sum_b \sigma_b^{-2}j_b^\top j_b.
\]
增加新独立束，新增项半正定，因此信息在 Loewner 顺序上不减；
当 $\mathcal I$ 正定时，逆矩阵所代表的局部方差界不增。
但若所有束都只约束平行墙法向，沿走廊方向的列可能接近零，
$\mathcal I$ 仍退化，更多同方向束不产生缺失的可观测性。
占据距离场在障碍内部为零、最近障碍切换处不可微，上式只限可微邻域；
不能把它当成全地图全局证明。

相关性也限制信息增长。设 $B$ 个束误差方差均为 $\sigma_e^2$，
两两相关系数为 $\rho\ge0$，则平均误差
$\bar e=B^{-1}\sum_be_b$ 满足
\begin{align}
\Var(\bar e)
&=\frac1{B^2}\left[B\sigma_e^2+B(B-1)\rho\sigma_e^2\right]\\
&=\frac{\sigma_e^2}{B}\,[1+(B-1)\rho].
\end{align}
等效独立束数
$B_{\rm eff}=B/[1+(B-1)\rho]$。
说明性 $B=120,\rho=0.1$ 得 $B_{\rm eff}\approx9.30$。
因此 GPU 允许用全部激光束，并不意味着可以按独立乘积无限提高置信度。
共享地图误差、外参误差和运动畸变会相关；应增加相关噪声模型或按结构抽样，
并用外部真值检验后验校准。

\subsection{多假设比单高斯局部定位更适合当前歧义}
重复走廊可能同时支持多个位姿模式。
把多峰后验压成单个均值可能落在两个合法模式之间的非法位置。
合理结构是维护模式权重 $\pi_a$、模式内粒子/局部网格及各自状态，
\[
b(q)=\sum_{a=1}^{A_m}\pi_a b_a(q),\qquad
\pi_a\ge0,\quad\sum_a\pi_a=1.
\]
并行更新所有模式后，再依据观测证据、动态一致性和任务风险决策。
若使用局部扫描优化作为粒子提议，优化后的点不再服从原运动提议，
必须计算提议修正、使用合适 MCMC 接受率，或明确改为混合启发式算法。
不能在每个粒子上做一次梯度下降，仍声称保持原 AMCL 后验。

这也是为什么不能简单用 GPU NDT/ICP 替换 AMCL：
单次局部配准通常给出一个局部最优位姿，无法自动提供原粒子滤波的多峰表示及全局覆盖。
它可以作为多假设框架内的局部精化器；更大批量有价值，但优化器名称不构成精度保证。

\section{扫描匹配：扩大搜索范围并提高局部拟合精度}
\base{localization.tex 的 Karto 相关匹配部分}
\subsection{候选位姿与激光点的二维批处理}
设候选位姿 $q_s=(p_s,\theta_s)$，$s=1,\dots,N_s$，
激光点 $y_b$ 共 $B_s$ 个，相关栅格为 $G_c$。
原评分可写为
\begin{equation}
c_s=\frac1{B_s}\sum_{b=1}^{B_s}
\ind_{\rm valid}(p_s+R(\theta_s)y_b)\,
G_c(p_s+R(\theta_s)y_b).
\end{equation}
越界项贡献零，分母仍为总点数 $B_s$。
再乘原稿中带下限的距离与航向惩罚 $\chi_d(q_s)\chi_\theta(q_s)$。
每个 $(s,b)$ 独立，随后按 $s$ 归约，再全局求最大，是直接的 GPU 结构。

若搜索平移半宽 $L_x,L_y$、角度半宽 $L_\theta$，
步长分别为 $r_x,r_y,r_\theta$，
候选数约为
\[
N_s=
\left(\left\lfloor\frac{2L_x}{r_x}\right\rfloor+1\right)
\left(\left\lfloor\frac{2L_y}{r_y}\right\rfloor+1\right)
\left(\left\lfloor\frac{2L_\theta}{r_\theta}\right\rfloor+1\right).
\]
三维分辨率均减半，候选数约增为 8 倍；
搜索半宽均加倍也约增为 8 倍。
这正是 CPU 容易限制但 GPU 可能承担的计算扩张。
收益分别是减少离散候选量化、增加初值误差容忍范围，不应混为“精度提高 8 倍”。

若真实解不在窗口内，细化窗口内分辨率不能找回它。
若窗口很大且地图重复，最大相关分数可能来自错误走廊；
算得更多还会增加误匹配候选。因此需同时使用先验一致性、多个高分模式及几何验证。

\subsection{保证不漏候选的分层搜索}
对一个候选子集合（搜索盒）$\mathcal B$，第 $b$ 个点可能落入的栅格集合为
$X_b(\mathcal B)=\{p+R(\theta)y_b:q\in\mathcal B\}$。
定义该点的上界
\[
U_b(\mathcal B)=\max_{x\in X_b(\mathcal B)}\widetilde G_c(x),
\]
其中 $\widetilde G_c$ 越界为零。
任何候选均满足逐点评分不超过 $U_b$，故
\[
c(q)\le\frac1{B_s}\sum_b U_b(\mathcal B).
\]
若 $0\le\chi_d,\chi_\theta\le1$，再令
$U_\chi(\mathcal B)\ge\max_{q\in\mathcal B}\chi_d(q)\chi_\theta(q)$，
得到
\begin{equation}
c(q)\chi_d(q)\chi_\theta(q)
\le U(\mathcal B)
=\frac{U_\chi(\mathcal B)}{B_s}\sum_bU_b(\mathcal B).
\label{eq:bnb}
\end{equation}
$U_\chi=1$ 总是可用但较松；也可按盒内到先验的最短距离求更紧界。
维护已得到分数 $S_{\rm best}$。
若 $U(\mathcal B)<S_{\rm best}$，该盒不可能包含同等或更好解，可删除；
否则分裂搜索盒，直至叶节点逐个计算。
严格小于的剪枝保留并列最佳解，便于复现原 Karto 的并列位置均值与圆均值。

最大值金字塔可快速查询 $U_b$；
旋转造成的点移动包围域必须覆盖整个搜索盒，不能只查盒中心。
这种有上界的分支定界在有限候选集上可保持最优值，
与“粗层只保留 top-1 再细化”的启发式不同。
后者可能丢掉粗层分数略差、细层正确的模式。
GPU 可批量展开一批搜索盒，CPU 管理不规则队列。
Cartographer 的作者论文提供了 scan-to-submap 分支定界先例\cite{cartographer}；
本文公式是针对原稿评分结构的适用条件分析，非现有 Karto CUDA 插件声明。

\subsection{连续精化：每束并行构造小型正规方程}
原稿局部平移搜索步长是 5 mm，不等于定位误差上界。
在离散候选附近，可换用可微插值场做连续局部优化，
消除“只能选择离散候选”造成的限制。
此时必须明确目标已改变，建议保持多个粗搜索模式各自精化。

使用统一量纲变量 $z=(p_x,p_y,\ell_\theta\theta)^\top$，
$\ell_\theta>0$ 为选定尺度长度。
令连续距离场 $\bar D$ 的单位为 m，
令 $\sigma_b>0$ 为第 $b$ 束的距离残差尺度（m）。标准化残差
$r_b(z)=\bar D(p+R(\theta)y_b)/\sigma_b$。
先定义二次目标
\begin{equation}
E(z)=\frac12\sum_b r_b(z)^2+
\frac12(z-z_0)^\top P_z^{-1}(z-z_0),
\label{eq:scanls}
\end{equation}
其中 $z_0$ 是指定先验均值，$P_z$ 是正定先验协方差，不与其他模块重复使用同一证据。
在当前点 $z$，
\[
r_b(z+\Delta z)\approx r_b(z)+j_b\Delta z,\quad
j_b=\frac{\nabla\bar D(e_b)^\top}{\sigma_b}
\begin{bmatrix}I_2&R(\theta)Jy_b/\ell_\theta\end{bmatrix}.
\]
代入目标，展开平方，对 $\Delta z$ 求导：
\begin{align}
\nabla_{\Delta z}E_{\rm lin}
&=\sum_bj_b^\top(r_b+j_b\Delta z)
+P_z^{-1}(z-z_0+\Delta z).
\end{align}
令导数为零并收集 $\Delta z$：
\begin{equation}
\underbrace{\left(\sum_bj_b^\top j_b+P_z^{-1}\right)}_{H_z}
\Delta z
=-\underbrace{\left(\sum_bj_b^\top r_b+P_z^{-1}(z-z_0)\right)}_{g_z}.
\label{eq:scangn}
\end{equation}
每个激光束独立贡献 $3\times3$ 对称矩阵的 6 个分量和向量的 3 个分量。
GPU 并行计算、归约这 9 个量；最终 $3\times3$ 线性求解本身极小。
只有一个起点、少量束时 CPU 可能更快；
多个假设、多分辨率、多个子地图同时评估才更适合 GPU。
若矩阵病态，可使用 LM 阻尼 $(H_z+\mu I)\Delta z=-g_z$；
变量尺度已统一，$\mu I$ 才有明确的数值意义。

对动态障碍离群点可把 $\frac12r^2$ 替换为 Huber 损失
\[
h_\kappa(r)=
\begin{cases}
r^2/2,&|r|\le\kappa,\\
\kappa|r|-\kappa^2/2,&|r|>\kappa.
\end{cases}
\]
其导数是 $r$ 或 $\kappa\operatorname{sgn}(r)$，
因此 IRLS 权重为
\[
a_b=\frac{h_\kappa'(r_b)}{r_b}
=\min(1,\kappa/|r_b|),\quad a_b=1\text{ 当 }r_b=0.
\]
固定权重的一轮更新只需将式 \eqref{eq:scangn} 中
$j_b^\top j_b,j_b^\top r_b$ 改为 $a_bj_b^\top j_b,a_bj_b^\top r_b$。
这是对新目标的迭代近似，不是原 Karto 算法的等价替换。

连续插值改善可优化性，同时改变目标表面。
最近障碍切换、未知区域、粗地图内的无信息平台都会限制梯度，
局部 Hessian 的逆也不自动是校准协方差。
需要步长/信赖域与更新后的目标下降检查；
多峰之间仍由粗搜索覆盖，不能靠一次局部求解保证全局正确。

\subsection{更高维模型与误匹配控制}
GPU 可让每个扫描候选同时估计微小时间偏移、外参扰动或扫描内运动参数。
如果额外变量 $\xi$ 改变端点 $e_b(q,\xi)$，只需将
$j_b$ 增加
$\nabla\bar D(e_b)^\top\partial e_b/\partial\xi$ 这些列。
代价是维度增加、信息矩阵更易退化；
必须有足够运动激励和先验约束，不能同时任意释放位置、外参和时间偏移。

扩大搜索还改变多重比较：若每个错误候选独立超过阈值概率为 $p_f$，
$N_s$ 个错误候选至少一次超过阈值概率为
$1-(1-p_f)^{N_s}$。
实际评分相关，独立模型只是解释；
无论如何不能在大幅扩展搜索后沿用未经复验的单候选接受阈值。
多假设保持、回环几何一致性与后续时序证据，比单纯选最大分更必要。

\section{位姿图、规划与地图更新：适合局部卸载的环节}
\subsection{位姿图的可并行部分及其规模门槛}
\base{localization.tex 的 Ceres 图优化部分}
对边 $(a,b)$，设
$d_{ab}=R_a^\top(p_b-p_a)$、观测 $\bar d_{ab},\bar\theta_{ab}$，
未白化残差
\[
e_{ab}=
\begin{bmatrix}d_{ab}-\bar d_{ab}\\
\wrap(\theta_b-\theta_a-\bar\theta_{ab})\end{bmatrix}.
\]
在角差非分支点，因 $(R_a^\top)'=-JR_a^\top$，雅可比为
\begin{equation}
A_{ab}=\frac{\partial e_{ab}}{\partial q_a}
=\begin{bmatrix}-R_a^\top&-Jd_{ab}\\0&-1\end{bmatrix},
\quad
B_{ab}=\frac{\partial e_{ab}}{\partial q_b}
=\begin{bmatrix}R_a^\top&0\\0&1\end{bmatrix}.
\end{equation}
其中最后一行前两项是零；两个矩阵均为 $3\times3$。
信息矩阵 $\Omega_{ab}=\Sigma_{ab}^{-1}$，白化因子应满足
$U_{ab}^\top U_{ab}=\Omega_{ab}$。
线性化后，一条边向全局矩阵贡献
\[
H_{aa}{+}=A^\top\Omega A,\quad
H_{ab}{+}=A^\top\Omega B,\quad
H_{bb}{+}=B^\top\Omega B,
\]
并向梯度贡献
$g_a{+}=A^\top\Omega e,\ g_b{+}=B^\top\Omega e$，
对称块 $H_{ba}=H_{ab}^\top$。
独立边评估适合 GPU，但同节点写入要分段归约或受控原子累加。

固定连通图锚点、统一变量尺度并加入正阻尼后，
求 $(H+\mu I)\Delta=-g$。
GPU 稀疏 Cholesky 或预条件共轭梯度（PCG）适合较大的稀疏系统。
PCG 的一个核心操作可以不显式形成 $H$：
\[
v\longmapsto J^\top(Jv)+\mu v.
\]
边局部矩阵乘法和节点归约为主要并行部分。
迭代数依赖条件数；正规方程还会平方 $J$ 的条件数，不能只按非零元数估计求解时间。
若图断连，单个锚点不能消除所有分量的规范自由度。

较新 Ceres 官方文档列出 CUDA 稠密求解、CGNR 及 cuDSS 稀疏后端选项\cite{ceresnew}；
本项目的实际版本与编译依赖需要另外确认，不能仅改配置名就假定 GPU 已启用。
二维中小位姿图的状态很少且优化非每个控制周期执行，
本报告不把全量图迁 GPU 排在 MPPI/匹配之前。

当前损失是 None。
用鲁棒边权、异常回环验证或增量更新可能比硬件移植更改善图质量，
但它们改变模型/求解流程。
Huber 只降低大残差影响，无法保证识别低残差却错误的重复结构回环；
需要跨扫描证据与约束一致性验证。

\subsection{Smac 搜索与更大航向修复族}
Smac 的优先队列弹出、动态分支与地图相关扩展较不规则。
GPU 可批量验证运动原语、计算启发式场，CPU 保留队列管理。
若同时弹出一批节点，应重新检查搜索最优性/停止条件；
不能将原有算法性质直接转移到无序 GPU 扩展。

项目的航向修复是固定位置序列上的分层 DAG：
\[
D_i(b)=\min_{a\in\Theta_{i-1}}
\{D_{i-1}(a)+c_i(a,b)\},
\]
其中非法边 $c_i=+\infty$，合法边为角行程，
每层当前至多 4 个候选。DP 标量计算很少，真正昂贵的是
每条边沿插值的填充碰撞及软代价比较。
可先 GPU 并行生成所有边可行性与代价，CPU 顺序 DP。

深度重构可把候选数从 $A=4$ 扩到 $A'$ 个方向，
边数从约 $16L$ 变为 $A'^2L$（$L$ 为段数）。
在候选族嵌套、原边代价与接纳规则不变且完整求解时，
原最优路径仍在扩大的图里，因此
\[
J^*_{\rm enlarged}\le J^*_{\rm original}.
\]
这是“更多计算不会变差”的明确有限图结论；
它只针对累计角行程，不保证整体到达时间或风险下降。
固定位置不变也无法解决原路径位置不合理的问题。
若增加横向位置候选，可进一步形成带宽受限的状态格优化，
但碰撞边数与空间规模会一起增加，需要配置空间缓存帮助摊销。

\subsection{射线、膨胀与一致性监测}
地图射线可按束并行遍历栅格，计算 pass/hit 整数统计，
GPU 原子加法在无溢出时保持整数和的顺序不变性；
必须复现越界、端点计数、未知区和严格阈值。
图优化改变历史扫描位姿后，累计图需要重新投影或维护可撤销贡献，
仅追加新射线会混用新旧几何。

膨胀可由距离场加指数核生成，但原 Nav2 的有限半径、
特殊值 253/254/255、层合并和未知空间策略必须逐项保留。
“到最近障碍的距离一样”不自动保证最终 costmap 每个字节相同。

定位质量监测最多 90 束、50000 个栅格段；
可合并到已有 GPU 地图服务，但单独启动小核收益不确定。
它的输出是几何一致性规则，增加射线数不能把它转成真值定位精度。
EKF 只有 15 维状态、四元数更新常数级，
更复杂 GPU 计算也不能创造偏航绝对观测或消除轮速系统偏差，
因此这类模块应继续放在 CPU 关键控制路径。

\section{固定耗时下的三类深度重构}
\label{sec:deep}
本节回应核心设计意图：\textbf{允许新算法比旧算法使用更多总运算和显存，
但让 GPU 将在线耗时压回原预算以内。}
以下不是对现有系统已经取得效果的描述，而是具备数学目标、预算变量和验证条件的研究方案。

\subsection{方案一：多模式、多场景的风险敏感采样控制}
\subsubsection{把不确定性纳入预测，而非只增加名义轨迹}
原 MPPI 对每个控制样本主要评估一条确定初始状态与既定运动模型下的轨迹。
定义不确定量
\[
\xi=(q_0,\tau_v,\kappa_{\rm slip},a_b,\alpha_b,\mathcal O_{0:H}),
\]
分别表示初始位姿、执行时间常数、侧滑参数、制动能力和未来障碍。
给定控制序列 $U_n$ 与场景 $\xi_m$，
\[
x_{n,m,\ell+1}=f(x_{n,m,\ell},u_{n,\ell},\xi_m),\qquad
J_{nm}=J(U_n;\xi_m).
\]
场景应来自经过校准的定位后验、动力学参数区间/分布及障碍预测，
不能把同一地图/标定误差对不同激光束独立抽样。
全局偏差应在一个场景内共享，才保留相关误差的物理含义。

计算张量从 $(K,H)$ 扩成 $(K,M,H)$，$n,m$ 相互独立。
同时对所有候选使用同一组 $\xi_m$，可形成 common random numbers：
比较两个候选时
\[
\Var(\widehat J_A-\widehat J_B)
=\Var(\widehat J_A)+\Var(\widehat J_B)
-2\Cov(\widehat J_A,\widehat J_B).
\]
如果共享场景使协方差为正，候选差值噪声下降。
这不保证所有问题都正相关，但避免了仅因抽到不同难度场景而错误排序。

\subsubsection{从平均风险到 CVaR 的逐步构造}
仅最小化均值
$\bar J_n=M^{-1}\sum_mJ_{nm}$
可能接受“多数很好、少数碰撞”的动作。
对风险置信水平 $\beta\in(0,1)$ 定义上尾 CVaR：
\begin{equation}
\CVaR_\beta(J)
=\min_{\eta\in\R}
\left[\eta+\frac1{1-\beta}\E(J-\eta)_+\right].
\label{eq:cvar}
\end{equation}
$\eta$ 与 $J$ 同单位，$(a)_+=\max(a,0)$。
对连续分布，在可微点
\[
\frac{d}{d\eta}\left[\eta+\frac1{1-\beta}\E(J-\eta)_+\right]
=1-\frac{\Prob(J>\eta)}{1-\beta}.
\]
令导数为零得 $\Prob(J>\eta)=1-\beta$，
故最优 $\eta$ 为 $\beta$ 分位点；代回得到上尾平均。
对有限样本则求
\[
\widehat{\CVaR}_\beta(J_n)
=\min_\eta\left[\eta+\frac1{(1-\beta)M}\sum_m(J_{nm}-\eta)_+\right],
\]
可通过排序或分位数选择精确求该经验目标，
边界样本按分数权重处理，不必硬取整数尾部数量。

新轨迹分数可定义
\begin{equation}
S_n^{\rm risk}=\bar J_n+\lambda_R\widehat{\CVaR}_\beta(J_n)+C_n^{\rm control},
\label{eq:riskcost}
\end{equation}
$\lambda_R\ge0$ 为同标度风险权重，$C_n^{\rm control}$ 是与采样方案匹配的控制修正。
然后仍可按第 4 节的熵正则推导对 $S_n^{\rm risk}$ 加权。
这改变了优化问题，目标从名义平均偏好转向尾部风险偏好。
CVaR 的标准来源见 Rockafellar--Uryasev\cite{cvarref}。
$M=16,\beta=0.9$ 时尾部只有约 $1.6$ 个样本的有效质量，
风险估计很粗；不能把很高的 $\beta$ 配很少场景当成可靠风险度量。

\subsubsection{有限场景零碰撞并非任意小概率保证}
对预先固定的某个动作，真实碰撞概率为 $p$，
$M$ 个独立同分布场景都不碰撞的概率为 $(1-p)^M$。
若用置信失误率 $\alpha$ 排除过大的 $p$，
令 $(1-p_{\rm upper})^M=\alpha$ 得
\[
p_{\rm upper}=1-\alpha^{1/M}.
\]
但控制器从 $K$ 个动作中选择，选择行为会利用场景样本。
若全部候选在抽场景前固定，按 union bound 为每个候选分配 $\alpha/K$，
可得同时有效的保守上界
\begin{equation}
p_{\rm upper}=1-(\alpha/K)^{1/M}.
\label{eq:scenario}
\end{equation}
不同候选的碰撞事件无需相互独立；需要每候选的场景抽样条件成立。
若候选依据同一场景数据继续自适应优化，则这个有限固定集合证明不再直接适用；
需独立验证场景或相应的统计学习保证。

因此 $K=1600,\alpha=0.05$ 时，
$M=32$ 的上界约 $0.277$，$M=256$ 也仍约 $0.0397$。
这些是保守统计界，说明“小批量零碰撞”不足以证明极低风险。
场景模型失配时即使统计界形式正确也不能保证实物风险；
对制动等关键量可额外采用物理可信的区间最坏情形。
GPU 能增加检查覆盖，不能消除模型校准这一条件。

\subsubsection{多模式输出：避免把左绕和右绕平均到障碍上}
可行控制集合通常非凸。
若左绕序列 $U_L$ 与右绕序列 $U_R$ 都安全，
其平均 $(U_L+U_R)/2$ 可能直行穿过障碍。
因此“每条样本都安全”不保证 MPPI 加权均值安全。

建议按停车、保持路径、左绕、右绕等行为族分别进行采样和权重更新，
得到各族候选 $\bar U_a$；
重新评估 $\bar U_a$ 本身及实际第一步过渡，
再从通过验证者中选择。
模式切换加入滞回，避免噪声导致左右来回切换。
重新评估是必要计算，不可因 GPU 已检查样本就省略。
这项重构改变的不仅是样本数，还改变了输出决策结构。

\subsubsection{在原预算内如何选轨迹数、场景数与步数}
基线 $KH=89600$。
一个\emph{预算候选}为 $K=800,M=8,H=56$、两轮更新，
状态展开数为 $716800$，约基线 8 倍。
如果式 \eqref{eq:expansion} 的实测允许扩张量至少为 8，
且新增风险归约、地图访问、最终候选复验也包含在时间中，
它才有机会满足同耗时要求。
这不是已经验证能实时的配置。

单轮同时保存 $K M H$ 个三维 float 位姿需
$800\times8\times56\times3\times4=4.3008$ MB。
若每场景维护更多状态和代价，内存相应增加。
可在轨迹展开时累积可加代价，减少全轨迹存储；
路径对齐等需要历史信息时则不能无条件流式丢弃。
若地图查询是主要瓶颈，新增场景未必按算力比例加速。
可采用“低成本全体筛选 + 高风险/高质量候选增加场景”的分层预算，
但筛选误差和自适应选择的统计偏差必须独立评估。

\subsection{方案二：显式位姿概率网格与层次多假设定位}
\subsubsection{用显存换分布覆盖，而非只保存有限粒子}
在离散状态空间
$\mathcal X=\{(i,j,k):0\le i<N_x,0\le j<N_y,0\le k<N_\theta\}$
上保存概率质量 $b_t(i,j,k)$，
$\sum_{i,j,k}b_t=1$。
给定离散运动转移矩阵 $P(x|x')$，
Bayes 递推为
\begin{align}
b_t^-(x)&=\sum_{x'}P(x|x')b_{t-1}^+(x'),\\
b_t^+(x)&=\frac{L_t(x)b_t^-(x)}{\sum_yL_t(y)b_t^-(y)}.
\label{eq:gridbayes}
\end{align}
这是选定有限状态模型的完整分布递推，无粒子重采样丢失模式问题；
但有网格量化、截断域和模型误差。
精确离散递推不等于精确连续定位。

若直接使用稠密转移矩阵，复杂度为状态数平方，显然不适合。
重构关键在利用运动核的局部性和按方向的平移不变结构。

\subsubsection{由全向运动模型得到卷积}
对旧方向层 $k'$，把平移噪声及平均里程计位移形成二维离散概率核
$G_{k'}(\Delta i,\Delta j)$。
它可包含随方向旋转的各向异性协方差与非零均值。
令角度增量核为 $A(k|k')$。
在“给定旧方向后平移与角增量独立、空间中运动模型平移不变”的假设下，
\[
P(i,j,k|i',j',k')
=A(k|k')G_{k'}(i-i',j-j').
\]
代入式 \eqref{eq:gridbayes}，交换有限求和：
\begin{align}
c_{k'}(i,j)&=
\sum_{i',j'}G_{k'}(i-i',j-j')b_{t-1}^+(i',j',k'),\\
b_t^-(i,j,k)&=\sum_{k'}A(k|k')c_{k'}(i,j).
\label{eq:convpredict}
\end{align}
第一步是每个方向的二维卷积，第二步是方向混合。
若角度核有效宽度为 $B_\theta$，直接方向混合工作
$O(GB_\theta N_\theta)$；平移卷积可局部 stencil 或 FFT。
FFT 用
$\operatorname{FFT}(G*b)=\operatorname{FFT}(G)\operatorname{FFT}(b)$，
每层工作约 $O(G\log G)$。
线性卷积必须零填充或保持足够 halo，
否则 FFT 的循环边界会把一侧概率错误送到另一侧。

若截断高斯运动核，必须记录丢弃尾质量并重新定义近似模型。
边界离开局部窗口的概率不能凭空归一化后假装没有发生；
可转入外层粗地图或保留“窗口外”质量。
测量更新可能放大先验近似误差，所以仅知道转移截断很小，
也不能无条件保证后验误差同样小。

\subsubsection{把全网格激光评分改写成相关运算}
当前 beam-skip 关闭且每束使用同一评分函数。
定义地图函数
\[
g(x)=\log\left[a\exp(-D_{\mathcal M}(x)^2/(2\sigma^2))+b_0\right].
\]
在固定方向 $\theta_k$、固定平移网格及一致的端点量化规则下，
令 $\Delta_{kb}$ 为第 $b$ 束对应的整数端点偏移。
创建稀疏扫描模板
$S_k(\Delta)=\sum_b\ind\{\Delta=\Delta_{kb}\}$；
重复落格的束以计数保留。
则该方向整幅位姿图的对数观测评分
\begin{align}
\ell_k(p)
&=\sum_bg(p+\Delta_{kb})\\
&=\sum_\Delta S_k(\Delta)g(p+\Delta).
\label{eq:senscorr}
\end{align}
第二行正是互相关，可用 FFT 实现。
在支持条件成立时，原 $O(N_\theta GB)$ 的点束评分，
转成约 $O(N_\theta G_{\rm pad}\log G_{\rm pad})$；
$G_{\rm pad}$ 是为线性相关扩大的图。
图的 FFT 可以复用，扫描模板随扫描变化。
这是一项比“增加线程”更深入的算法结构变化。

限制也明确：端点量化必须与原 grid 坐标一致；
逐粒子束跳过、位置相关束权、变化的可见性模型可能破坏简单相关形式；
越界必须用原截断距离对应的常数评分填充，
不能默认零或周期地图。
数值 FFT 会有误差，不能仅凭代数恒等式宣称逐位相同。
模板旋转的量化若改变端点，需单独分析。
从成本角度，$B=120$ 与 $\log G$ 的比较只是渐近判断，
常数、填充和 FFT 批量大小仍需测量。

\subsubsection{规模算例及合理的层次结构}
\begin{center}\small
\begin{tabular}{lrrr}
\toprule 位姿网格 & 状态数 & 单 float 场内存 & 直接 120 束评分量\\\midrule
局部 $4$ m，$4$ cm，64 方向 & $640000$ & $2.56$ MB & $76800000$\\
局部 $4$ m，$1$ cm，64 方向 & $10240000$ & $40.96$ MB & $1228800000$\\
全局 $20$ m，$1$ cm，64 方向 & $256000000$ & $1024$ MB & $30720000000$\\
\bottomrule
\end{tabular}
\end{center}
内存仅为一个标量场，未计预测/后验双缓冲、对数似然和 FFT 复数临时区。
最后一行说明显存能容纳不代表扫描周期能完成计算。
因此不建议第一步就做全地图毫米级三维稠密 Bayes。

更合理路线：全局较粗分布维持所有有意义模式，
对若干高质量且空间分离的模式建立局部细网格，
局部用连续匹配或粒子更新。
粗细层应是互斥分区的概率质量，或显式混合：
局部层权重由粗层分配，不能把同一状态质量在两层重复加和。
细化区域扩张、合并和删除时保持总质量；
不应只保留 top-1 后将其权重硬设为 1。
GPU 的额外存储在这里用于\emph{保留歧义}，
不是制造一个数值上很尖却可能错误的后验。

\subsection{方案三：预计算配置空间支撑更大动作图及主动消歧}
\subsubsection{空间换时间，让在线搜索能覆盖更多行为}
式 \eqref{eq:cspace} 把地图与足迹组合编译成位置--方向可行表。
当 $Q$ 次在线查询共享同一地图，原暴力成本约
$Q\,c\,N_oN_v$，
新成本约
$T_{\rm build}+Q\,t_{\rm lookup}+Q_{\rm near}t_{\rm exact}$。
需同时满足
\begin{equation}
T_{\rm build}+Q\,t_{\rm lookup}+Q_{\rm near}t_{\rm exact}
<Q\,cN_oN_v
\end{equation}
以及单周期更新峰值预算。
节省下的查询时间可用于增加航向/横移候选、多个全局通路、
更多控制情景或更细时间步长。
“是否值得预计算”取决于共享查询数 $Q$，
而不取决于单张表能否放进显存。

静态图和动态障碍可分层：静态配置空间预计算，
动态局部障碍每周期更新较小区域，
在线碰撞是两层碰撞结果的 OR。
若风险情景包含不同障碍轨迹，可以共享静态层，
按情景建立动态层或对动态障碍精查。
不得因为静态地图通过就忽略扫描中新出现的障碍。

\subsubsection{主动感知：让控制帮助定位解决歧义}
如果定位后验同时有多个模式，仅做更多计算可能仍无法区分重复走廊。
机器人可以选择能看到岔口/非对称结构的微小动作，
用下一次真实观测获得新信息。
这时额外 GPU 预算用于“候选动作 $\times$ 位姿假设 $\times$ 预测观测”。

设离散模式随机变量为 $A$，权重 $\pi_a$；
候选动作 $u$ 后的未来激光特征为 $Z$，模型密度 $p(z|a,u)$。
混合观测密度
$p(z|u)=\sum_a\pi_ap(z|a,u)$。
由 Bayes，
$\pi_a(z,u)=\pi_ap(z|a,u)/p(z|u)$。
先验熵 $H(A)=-\sum_a\pi_a\log\pi_a$，
期望后验熵为 $\int p(z|u)[-\sum_a\pi_a(z,u)\log\pi_a(z,u)]\,dz$。
相减并代入 Bayes，化简得
\begin{equation}
I(A;Z|u)=
\sum_a\pi_a\int p(z|a,u)
\log\frac{p(z|a,u)}{p(z|u)}\,dz.
\label{eq:infogain}
\end{equation}
可以把 $-\lambda_I I(A;Z|u)$ 加到任务代价，同时约束运动风险与任务进度。

直观地，两模式若预测同一测量高斯
$\mathcal N(\mu_a,\sigma_z^2)$，
其 KL 散度为
\[
D_{\rm KL}(\mathcal N(\mu_1,\sigma_z^2)\|
\mathcal N(\mu_2,\sigma_z^2))
=\frac{(\mu_1-\mu_2)^2}{2\sigma_z^2}.
\]
证明是将两正态对数密度相减，
对第一分布取期望：
$\E[(Z-\mu_2)^2-(Z-\mu_1)^2]=(\mu_1-\mu_2)^2$。
这只是两模式可区分度，不等于式 \eqref{eq:infogain} 的混合互信息；
但说明动作若扩大预测观测差异，有潜力增加消歧信息。
若所有可执行动作下预测都相同，任何计算能力也无法产生信息。

该方向最能体现“用算力提升算法能力”，同时开发与建模成本最高。
应放在可靠多假设定位和场景控制之后，
先验证预测观测模型，再引入主动动作。
同一激光既用于定位又用于场景生成时，应使用一个联合后验，
不能把重复使用的证据当作独立观测再次相乘。

\section{实现架构、数值精度与模型边界}
\subsection{建议的 CPU/GPU 分工}
\begin{figure}[htbp]
\centering
\begin{tikzpicture}[
box/.style={draw=accent,fill=lightaccent,rounded corners,align=center,text width=40mm,minimum height=17mm,font=\small},
arr/.style={-{Latex[length=2mm]},thick,draw=accent},
node distance=12mm and 9mm]
\node[box] (cpu) {CPU：时间戳、TF、快照\\状态机与预算管理};
\node[box,right=of cpu] (maps) {GPU：版本化常驻地图\\距离场、配置空间与概率场};
\node[box,right=of maps] (batch) {GPU：候选与情景批评估\\定位、轨迹、归约};
\node[box,below=of batch] (out) {CPU：结果时效与模式接纳\\最终指令及故障处理};
\node[box,left=of out] (verify) {边界精查与实际候选复验\\独立保护路径};
\draw[arr] (cpu)--(maps);
\draw[arr] (maps)--(batch);
\draw[arr] (batch)--(out);
\draw[arr] (out)--(verify);
\end{tikzpicture}
\caption{建议的职责边界。图中只描述数据与验证关系，不声明 GPU 具有硬实时保证。最终实际执行仍须通过防护。}
\end{figure}
常驻数组优先使用 structure-of-arrays：连续存储各状态分量，
按线程访问顺序选择轨迹/时间主序。
代价图随机读取可能受带宽与缓存支配，
而不是单纯受浮点运算峰值限制。
多个极小核可融合，固定依赖图可考虑 CUDA Graphs；
但动态候选规模、重试、分支与地图更新不能为捕获方便而改变语义。
CPU 构造、内存锁和设备到主机结果回传都要计时。

输入快照必须在异步上传完成前保持有效。
不能释放 ROS 消息或修改共用数组后，仍让 GPU 读取旧指针。
所有结果带 $(\text{地图版本},\text{扫描时间},\text{姿态时间},\text{算法参数版本})$；
过期结果按已定义策略处置。
当前几何实现是同步等待流完成，不能把它描述成已经具备异步截止时限保护。

\subsection{FP32、FP64 与“更快但改变决策”的风险}
设浮点单位舍入误差为 $u$。
对 $n$ 项普通顺序求和，在 $nu<1$ 且标准模型适用时，
\[
|\operatorname{fl}(\sum_i x_i)-\sum_ix_i|
\le\gamma_{n-1}\sum_i|x_i|,\qquad
\gamma_n=\frac{nu}{1-nu}.
\]
平衡树归约的每个分量路径深度约 $\lceil\log_2n\rceil$，
可用相应深度的 $\gamma$ 估计误差；
但这只覆盖加法，不覆盖地图索引、指数和三角函数误差。
因此 CPU/GPU 不逐位相同并不必然是算法错误，
同时“总体误差小”也不能豁免边界分支检查。

建议把连续评分与硬几何分开：
\begin{itemize}
\item MPPI 轨迹与评分可先评估 FP32 的误差，并对归约采用稳定顺序/更高精度累加；
其影响用式 \eqref{eq:softbound} 及实际命令差评估。
\item AMCL 用 log-sum-exp，检查地图取格、有效束一致与后验模式稳定性。
\item 硬碰撞边界、未知区、非有限输入继续使用保守界和必要的 FP64/CPU 精查。
\item 坐标先减局部原点能减小大数相减的绝对误差，但必须在更高精度中完成该减法，
不能把两个大世界坐标先降为 FP32 再指望相减恢复丢失信息。
\end{itemize}
RTX、数据中心 GPU 与 Jetson 的 FP64 吞吐、带宽和共享内存体系不同，
本报告不把某型号的理论峰值代入实际导航加速比。
小型矩阵或标量控制没有自动使用 Tensor Core 的优势；
保持几何毫米量级语义时，未经分析的低精度矩阵化尤其不合适。

\subsection{概率场景的坐标一致与相关性}
若初始位姿和地图障碍同时受同一个全局坐标变换影响，
相对几何可保持不变。
例如共同施加刚体变换 $T$ 后
$\dist(T\mathcal F,T O)=\dist(\mathcal F,O)$。
若只在场景中扰动机器人、却把本来同样由该 TF 变换得到的障碍固定，
可能虚增相对不确定性；反之也可能漏掉地图到局部坐标的相对误差。
场景应抽样物理相对误差源，再一致地传播到机器人与障碍，
不是把所有输出协方差当独立高斯相加。

若确有局部高斯位姿误差 $e_q\sim\mathcal N(0,P_q)$，
足迹点扰动的一阶近似为
\[
e_a=J_ae_q,\quad
J_a=\begin{bmatrix}I_2&R(\theta)Ja\end{bmatrix},\quad
P_a=J_aP_qJ_a^\top.
\]
沿障碍单位法向 $n$ 的投影标准差是
$\sigma_n=\sqrt{n^\top P_an}$。
在小角度、单局部法向、高斯且协方差校准条件下，
要求名义间隙
$d_a\ge\delta+\Phi^{-1}(1-\alpha_a)\sigma_n$
可限制该一维事件的概率；$\Phi$ 为标准正态分布函数，$\alpha_a$ 是分配给该事件的概率上限。
多个时刻/障碍的联合风险可用
$\sum_a\alpha_a\le\alpha_{\rm total}$ 作保守分配，
但需有效覆盖全部所考虑事件。
该条件不适用于未校准、多峰重定位或未观测障碍；
更大的 GPU 只能更充分计算模型，不能把模型假设变成事实。

\subsection{实时性与最终保护}
CUDA 高吞吐不等于可证明的最坏响应时间。
GPU 同时被桌面、其他推理或 SLAM 使用时可能排队，
首轮上下文初始化、扩容、页面迁移都会产生尾部耗时。
预分配、预热、独立流可以降低常见开销，但不构成硬实时证书。
最终保护与失效停车应具有独立明确的输入时效规则。

若重构后的算法超期，不能无条件继续用旧控制。
需要按当前速度、当前可用障碍和已验证停止策略决定动作。
“发零速度”仍只是停车请求，不是瞬间消除动量。
报告中的所有碰撞保证均针对相应的障碍表示和运动模型，不超出原稿明确的物理边界。

\section{如何验证“算法效果提高且耗时不增加”}
\subsection{四组实验，分离硬件与算法收益}
\begin{center}\small
\begin{tabular}{p{0.12\textwidth}p{0.22\textwidth}p{0.21\textwidth}p{0.34\textwidth}}
\toprule 组别 & 算法 & 后端 & 目的\\\midrule
A & 原算法、原规模 & 优化且公平的 CPU & 冻结基线\\
B & 原算法、原规模 & CUDA & 测移植等价性与计算提速\\
C & 重构算法、扩大问题 & CPU，可离线 & 识别算法本身收益与计算代价\\
D & 同 C & CUDA，固定时延预算 & 检验用户要求的最终组合\\
\bottomrule
\end{tabular}
\end{center}
C 组离线结果只用于判断算法能力，不能拿减速仿真或暂停时钟掩盖实时超期。
A 与 D 的闭环比较必须使用相同墙钟约束、传感输入、场景难度及全部尝试分母。
先冻结候选规模选择规则，再在未用于调参的场景上比较。
若 D 使用更好的模型、更大样本、更多场景，应做逐项消融确定贡献。

\subsection{应记录的时间、规模与质量}
\begin{longtable}{p{0.19\textwidth}p{0.71\textwidth}}
\toprule 项目 & 必须记录的内容\\\midrule\endfirsthead
\toprule 项目 & 必须记录的内容\\\midrule\endhead
端到端时延 & 输入快照建立到可执行命令；P50/P95/P99、最大观测值、超期次数；首次启动与稳态分开\\
细分成本 & CPU 准备/锁、上传、核、下载、同步、地图更新、边界复核、最终候选复验\\
实际负载 & $K,H,M,N_p,B,N_o,N_q$；MPPI 调用比例；CUDA 真调用比例；每幅地图复用次数\\
时效 & 扫描、TF、地图、后验与命令年龄；是否发生丢帧、排队或旧快照复用\\
资源 & 峰值显存、GPU 时钟/功耗状态、CPU 线程数、并发其他任务、编译选项\\
定位 & 对独立真值的平移/角误差、丢定位率、错误模式率、重定位时间、后验校准与配对覆盖率\\
控制 & 全部任务成功率、完成时间、最小实物间隙、碰撞/近碰撞、跟踪误差、停止次数与指令变化率\\
风险 & 预测风险与实际发生率是否匹配；动态障碍、执行延迟、侧滑、制动能力变化的分层结果\\
\bottomrule
\end{longtable}
不将 4 mm 目标容差当作定位精度。
不将成功子集误差改善解释为整体表现改善而忽略失败分母。
不将高 GPU 利用率、kernel 时间或有限零碰撞样本替代闭环指标。

\subsection{场景与消融如何针对本报告的假设}
应覆盖：宽阔静态区、窄道与旋转角点、早停高碰撞密度、稀疏无碰撞大批量、
重复走廊多峰定位、突然全局定位偏移、动态横穿、弱纹理几何/长走廊退化、
地图更新密集、输入过期与 GPU 故障。
场景需要同时含“GPU 合适”和“CPU 早停占优”的负载，防止选择性报告。

算法消融建议：
\begin{enumerate}
\item 原几何 vs 空间索引 vs 距离场筛选 vs 配置空间，比较相同碰撞语义及保守拒绝率。
\item 同预算下只增 $K$、只增 $M$、更好动力学、结构化控制、多模式输出分别比较。
\item 同预算下更多粒子、更多束、更多独立模式、概率网格/相关评分分别比较。
\item 扩大搜索窗口、减小格距、连续精化分别比较，记录错误匹配而非只看拟合残差。
\item 固定地图更新年龄 vs 双缓冲允许年龄，量化缓存收益是否由过期地图换来。
\end{enumerate}
若同一任务重复随机运行，使用成对种子和相同输入扰动；
对 MPPI 噪声复用的基线，先固定样本以隔离后端差异，
再另做多随机种子评价。
统计不确定性应按任务/运行聚类，不能把同一条轨迹的每一时刻当独立试验。

\subsection{本次实际完成的核验}
本次未运行新导航实验或新 CUDA 基准。
报告源文件中的 \code{verify_math.py} 使用固定种子独立核验：
MPPI 熵正则驻点与权重误差界、Omni 前缀展开、
恒定旋量积分、激光/位姿图雅可比的有限差分、
距离场保守下界、方向量化位移界、
有限网格 Bayes 转移分解、FFT 相关与直接束求和、
CVaR 经验形式，以及预算/显存算例。
14 类检查全部通过，结果写入 \code{math_validation.json}。扫描雅可比有限差分最大绝对差约 $6.18\times10^{-8}$，位姿图雅可比约 $5.72\times10^{-10}$，FFT 相关与直接求和约 $7.11\times10^{-15}$；均在脚本预设容差内。运行采用独立 CPU 容器、Python 3.10 与 NumPy 1.26.4，未连接机器人或 GPU。
这些检查支持公式与数值算例的自洽，不代表新算法在机器人上可实时或更安全。
完整证明仍以文中假设与推导为准，随机数值检查不是形式化证明。

\section{建议的研究次序与可审查的成功条件}
\subsection{阶段一：拿到真实预算，并建立可比较基线}
记录整条控制路径时延与分支占比；确认实际瓶颈是 MPPI、AMCL、扫描匹配还是几何/地图更新。
先做 B 组同算法 CUDA 评估，并与带空间索引/向量化的 CPU 对照。
通过条件是：语义差异全部解释，GPU 真执行，端到端 P99 与过期输入情况可量化。
没有这一步，无法为计算扩张系数 $\alpha$ 给出可靠上限。

\subsection{阶段二：先扩大模型覆盖，再按收益分配样本}
\textbf{导航控制路线：}
在 MPPI 经常使用或尾延迟占主导时，
先实现设备常驻地图与批评分，
然后比较“增加样本”与“增加情景 + 更真实动力学 + 模式分离”。
优先增加能解释现有失败原因的维度，而不是默认把全部预算乘到 $K$ 上。
保留原直接跟踪分支作为明确基线；
若新研究要替换该分支，也单独计入调用次数与行为变化。

\textbf{定位路线：}
固定地图优先稳定 AMCL 评分与多模式保持；
歧义明显时试验粗全局/细局部分布，
只有足够复用和地图规模时再启用 FFT 全网格评分。
在线 SLAM 优先扩大有验证的相关搜索，再增加多起点连续精化，
最后才考虑大图 GPU 求解。
任何“更低匹配残差”都需独立真值检验，防止过拟合错误地图位置。

\subsection{阶段三：建立共享预计算，让更大问题变得便宜}
以静态距离场、方向配置空间和版本化动态层作为共享基础设施，
让新增航向、横移、控制情景和定位模式能够复用相同地图计算。
近边界精确查询继续保留，使粗表保守性不封锁窄道。
成功条件是\emph{包含地图构建和边界复验}后仍满足预算，而不只是查表核很快。

\subsection{最终应接受怎样的结论}
可以接受的结论是：
“在冻结版本、给定设备与规定场景分布上，
新方案相对 A 组在相同或更低端到端时延及规定超期约束下，
改善预先指定的定位/导航指标，且未恶化约定的风险和可靠性约束。”
应同时公开失败场景与未获益负载。

目前最有根据的研究主张是：
\textbf{GPU 使原来代价过高的多模式、多情景、较细配置空间搜索成为可考察方案；
真正提高效果的原因是表示、模型与搜索覆盖变得更充分，
而硬件提供了承担额外计算的可能。}
是否确实达到“效果更好且耗时不增加”，需要上述固定预算成对实验回答，
不能从显存容量、运算峰值或原稿的几何零差异测试直接推出。

\appendix
\section{原稿定位索引与复现说明}
\begin{longtable}{p{0.31\textwidth}p{0.59\textwidth}}
\toprule 原稿文件 & 本报告使用的内容\\\midrule\endfirsthead
\toprule 原稿文件 & 本报告使用的内容\\\midrule\endhead
\code{introduction.tex} & 系统模式、CUDA 现有范围与工作区版本边界\\
\code{notation.tex} & 坐标方向、底盘速度定义、填充足迹与圆盘障碍\\
\code{localization.tex} & 实际 15 维 EKF、AMCL 束乘积/KLD、Karto、图残差与质量监测\\
\code{geometry.tex} & 碰撞谓词、容差、采样界、CUDA 分块与门限\\
\code{control.tex} & 混合分支、56 步/1600 轨迹、噪声复用、原地滤波、33 条制动候选\\
\code{native_costs.tex} & 七项原生 critic、整批索引依赖、原生轮廓与填充碰撞的区别\\
\code{evidence.tex} & 历史几何差分、早停负载反向加速比、历史运动实验不能外推\\
\code{discussion.tex} & 固定间隙与概率不确定性、物理安全边界\\
\code{appendix.tex} & 有效参数、实现入口与运行平台边界\\
\code{source_manifest.json} & 本次核对的 101 个文件及 SHA-256\\
\bottomrule
\end{longtable}
本报告源目录含各节 \LaTeX{}、构建脚本、数学检查脚本、
输入与依赖来源清单以及验证结果。
使用 XeLaTeX 与 Noto CJK/URW 字体可重建 PDF；
数值检查需要 Python 3 与 NumPy。
原始压缩包未被改写，项目源文件和配置未被本任务修改。

\begin{thebibliography}{99}
\small
\bibitem{cuda}
NVIDIA. CUDA C++ Best Practices Guide. 官方文档，访问 2026-10-15。
用于数据驻留、传输成本与性能分析原则，不作为本项目实测加速数据。
\url{https://docs.nvidia.com/cuda/cuda-c-best-practices-guide/index.html}.
\bibitem{edt}
Felzenszwalb P F, Huttenlocher D P. Distance Transforms of Sampled Functions.
Theory of Computing, 2012, 8(19):415--428.
\url{https://theoryofcomputing.org/articles/v008a019/}.
\bibitem{pba}
Cao T T, Tang K, Mohamed A, Tan T S.
Parallel Banding Algorithm to Compute Exact Distance Transform with the GPU.
I3D 2010:83--90. 作者项目页及 PBA+ 说明，访问 2026-10-15。
\url{https://www.comp.nus.edu.sg/~tants/pba.html}.
\bibitem{mppisrc}
Nav2 Maintainers. MPPI optimizer, version 1.1.20. 官方版本化源码。
本报告具体 critic 语义以用户原稿附录及其固定版本为准。
\url{https://github.com/ros-navigation/navigation2/blob/1.1.20/nav2_mppi_controller/src/optimizer.cpp}.
\bibitem{itmppi}
Williams G, Drews P, Goldfain B, Rehg J M, Theodorou E A.
Information Theoretic Model Predictive Control: Theory and Applications to Autonomous Driving.
arXiv:1707.02342, 2017.
\url{https://arxiv.org/abs/1707.02342}.
\bibitem{amclsrc}
Nav2 Maintainers. LikelihoodFieldModelProb, version 1.1.20. 官方版本化源码，访问 2026-10-15。
\url{https://github.com/ros-navigation/navigation2/blob/1.1.20/nav2_amcl/src/sensors/laser/likelihood_field_model_prob.cpp}.
\bibitem{cartographer}
Hess W, Kohler D, Rapp H, Andor D.
Real-Time Loop Closure in 2D LIDAR SLAM. ICRA 2016.
作者机构页面，访问 2026-10-15。
\url{https://research.google/pubs/real-time-loop-closure-in-2d-lidar-slam/}.
\bibitem{ceresnew}
Ceres Solver Maintainers. Solving Non-linear Least Squares.
官方文档，访问 2026-10-15。
该文档描述的 GPU 后端不能自动视为本项目已安装依赖的能力。
\url{https://ceres-solver.readthedocs.io/latest/nnls_solving.html}.
\bibitem{cvarref}
Rockafellar R T, Uryasev S. Optimization of Conditional Value-at-Risk.
Journal of Risk, 2000, 2(3):21--41.
作者提供的论文预印本（1999）及正式发表信息，访问 2026-10-15。
\url{https://sites.math.washington.edu/~rtr/papers/rtr179-CVaR1.pdf}.
\end{thebibliography}

\end{document}
